\chapter{振动和波}

\section{简谐振动}

\begin{problem}{3.1.1}
一平台在竖直方向做简谐振动, 振幅为 $5 \mathrm{~cm}$ , 频率为 $\frac{10}{\pi} \mathrm{s}^{-1}$ , 在平台到达最低点时, 将一木块轻轻地放在平台上 (设木块质量远小于平台质量. 问:

(1) 木块于何处离开平台？

(2) 木块能达到的高度比平台能达到的高度高多少?
\end{problem}
\clearpage


\begin{problem}{3.1.2}
图 3.1 中质量为 2kg 的物体从平衡位置拉开并释放后以 $\frac{\pi}{6}s$ 的周期在光滑的水平面上沿直线振动.

(1) 将它由平衡位置拉开 2cm 需要多大作

用力？

(2) 如果将一小质量块放在 2kg 质量块上, 它们之间的静摩擦因数为 0.1, 问: 2kg 物体的振动振幅不超过何值时, 小质量块对它无相对运动? (假定附加小质量块后对振动周期的影响可以忽略)
\begin{center}
\begin{tikzpicture}[>=Stealth, scale=1]
  % Hatched walls and floor
  \fill[pattern=north east lines] (-0.22, -0.22) rectangle (0, 1.35);
  \fill[pattern=north east lines] (0, -0.22) rectangle (5.2, 0);
  \fill[pattern=north east lines] (5.2, -0.22) rectangle (5.42, 1.35);
  \draw[thick] (0, 1.35) -- (0, 0) -- (5.2, 0) -- (5.2, 1.35);

  % Mass block
  \draw[thick, fill=white] (2.15, 0) rectangle (3.05, 0.90);
  \node at (2.60, 0.45) {$2\,\mathrm{kg}$};

  % Left and right springs
  \draw[thick, decorate, decoration={zigzag, segment length=5.5, amplitude=3.2, pre length=14, post length=14}] (0, 0.45) -- (2.15, 0.45);
  \draw[thick, decorate, decoration={zigzag, segment length=5.5, amplitude=3.2, pre length=14, post length=14}] (3.05, 0.45) -- (5.2, 0.45);
\end{tikzpicture}

\end{center}
\end{problem}
\clearpage


\begin{problem}{3.1.3}
在劲度系数为 $k$ 的弹簧下悬挂一质量为 $m$ 的盘子, 处于静止状态. 有一质量也为 $m$ 的物体, 自比原静止的盘子高 $h$ 处由静止落下, 与盘子做完全弹性碰撞 (碰撞时机械能不损失), 求:

(1) 头两次碰撞期间盘子的运动；

(2) 用作图法求第三次碰撞的时间(取第一次碰撞时 t=0).
\end{problem}
\clearpage


\begin{problem}{3.1.4}
两个完全相同的圆柱体, 它们的轴平行, 且在同一水平面上, 相距为 $2L$ , 以相同的角速率按图 3.3 所示方向绕轴快速转动, 在圆柱体上放一均质木板, 木板与圆柱体之间的滑动摩擦因数为 $\mu$ . 开始把木板放在平衡位置偏右 $x_0$ 的位置, 且给它一个向右的初速度 $v_0$ , 求木板的运动.
\begin{center}
\begin{tikzpicture}[>=Stealth, scale=1]
  % Two rotating cylinders
  \draw[thick] (-1.45, 0) circle (0.95);
  \draw[thick] (1.45, 0) circle (0.95);

  % Rotation arrows (left clockwise, right counterclockwise)
  \draw[->, thick] (-1.45, 0) ++(155:0.72) arc (155:100:0.72);
  \draw[->, thick] (1.45, 0) ++(25:0.72) arc (25:80:0.72);

  % Horizontal plank resting on top
  \draw[thick, fill=white] (-2.75, 0.95) rectangle (2.75, 1.13);

  % Distance 2L between cylinder axes
  \draw[thin] (-1.45, -1.08) -- (-1.45, -1.52);
  \draw[thin] (1.45, -1.08) -- (1.45, -1.52);
  \draw[<->, thin] (-1.45, -1.32) -- (1.45, -1.32) node[midway, above=2pt] {$2L$};
\end{tikzpicture}

\end{center}
\end{problem}
\clearpage


\begin{problem}{3.1.5}
质量为 $m$ 和 $2m$ 的两质点系在一个不可伸长的轻绳两端，并搭在光滑的滑轮上，在 $m$ 的下端用原长为 $a$ 、劲度系数 $k = \frac{mg}{a}$ 的质量可忽略的弹簧与另一个质量为 $m$ 的质点相连。

开始时从下列状态:三个质点均静止;绳被拉紧;弹簧为原长开始运动.求:

(1) 任何时刻两质量为 m 的质点间的距离；

(2) 绳子张力；

(3) 弹簧对质点的作用力.
\end{problem}
\clearpage


\begin{problem}{3.1.6}
一质点被一根轻弹性绳系于固定点 $A$ ，绳的固有长度为 $a$ ，当悬挂质点平衡后，绳的长度为 $a + b$ ，今质点从 $A$ 点静止下降，求质点落到最低点与 $A$ 点间的距离以及此过程所用的时间。
\end{problem}
\clearpage


\begin{problem}{3.1.7}
滑轮上绕有不可伸长的轻绳, 绳上悬一弹簧, 弹簧的另一端挂一质量为 $m$ 的物体, 当滑轮等角速度转动时, 物体以等速度 $v_{0}$ 下降, 此时弹簧伸长量为 $\lambda$ . 如将滑轮突然停住, 求此后运动中弹簧的最大伸长量和绳子的最大张力.
\begin{center}
\begin{tikzpicture}[>=Stealth, scale=1]
  % Pulley
  \draw[thick] (-0.72, 2.80) circle (0.72);

  % String and spring
  \draw[thick] (0, 2.80) -- (0, 2.25);
  \draw[thick, decorate, decoration={zigzag, segment length=5.5, amplitude=3.2}] (0, 2.25) -- (0, 1.15);
  \draw[thick] (0, 1.15) -- (0, 0.72);

  % Mass block and initial velocity
  \draw[thick, fill=white] (-0.42, 0.08) rectangle (0.42, 0.72);
  \node at (0, 0.40) {$m$};
  \draw[->, thick] (0, 0.08) -- (0, -0.58) node[below=1pt] {$v_0$};

  % Elongation lambda and t=0 mark
  \draw[thin] (-0.08, 1.15) -- (-0.82, 1.15);
  \draw[thin] (-0.42, 0.40) -- (-0.82, 0.40);
  \draw[<->, thin] (-0.64, 0.40) -- (-0.64, 1.15) node[midway, left=2pt] {$\lambda$};
  \node[anchor=north east] at (-0.46, 0.34) {$t=0\text{ 时}$};

  % Coordinate axis x
  \draw[thin] (0.72, 0.40) -- (1.12, 0.40) node[right=2pt] {$0$};
  \draw[->, thick] (0.92, 0.40) -- (0.92, -1.05) node[below=1pt] {$x$};
\end{tikzpicture}

\end{center}
\end{problem}
\clearpage


\begin{problem}{3.1.8}
如果通过实心的均质分布的无自转动的地球沿一条直径钻一小隧道，一块石头从地球的一边从静止开始让它掉入隧道，经 $t_{1}=\frac{\pi}{\omega_{0}}$ 时间后在地球的另一边会看到它。如果石头不是掉入而是以速度 $v_{0}$ 抛入，问 $v_{0}$ 应当多大才能使它经 $t_{2}=\frac{1}{2}t_{1}$ 时间后在地球的另一边出现？答案用 $\omega_{0}$ 和地球的半径R给出。
\end{problem}
\clearpage


\begin{problem}{3.1.9}
一根不计质量、不可伸长的绳子，穿过光滑的水平桌面上的小孔，两端分别连着质量分别为 $m_{1}$ 和 $m_{2}$ 的两个物体。 $m_{1}$ 在水平桌面上做半径为 $r_{0}$ 的圆周运动时， $m_{2}$ 静止不动，如图3.8所示。今给 $m_{2}$ 一个向下的极小的冲量，证明 $m_{2}$ 将做简谐振动，并求出其振动周期（小孔的摩擦力可忽略不计）。
\begin{center}
\begin{tikzpicture}[>=Stealth, scale=1]
  % Horizontal tabletop in perspective
  \draw[thick] (-2.4, -0.70) -- (1.6, -0.70) -- (2.4, 0.70) -- (-1.6, 0.70) -- cycle;

  % Dashed circular orbit
  \draw[dashed, thick] (0, 0) ellipse (1.25 and 0.46);

  % Strings and masses
  \coordinate (M1) at ({1.25*cos(28)}, {0.46*sin(28)});
  \draw[thick] (0, 0) -- (M1);
  \draw[thick] (0, 0) -- (0, -2.15);

  % Small hole at center
  \draw[thick, fill=white] (0, 0) circle (0.065);

  % Masses m1 and m2
  \fill (M1) circle (0.075) node[right=3pt] {$m_1$};
  \fill (0, -2.15) circle (0.075) node[right=3pt] {$m_2$};
\end{tikzpicture}

\end{center}
\end{problem}
\clearpage


\begin{problem}{3.1.10}
一质量为 $M$ 的质点在势能 $V(x) = -kxe^{-ax}(k$ 和 $a$ 均为正值、常量)的影响下沿 $x$ 轴运动, 求平衡位置和在此平衡位置附近做小幅振动的周期. 若 $k$ 和 (或) $a$ 是负的常数, 情况如何?
\end{problem}
\clearpage


\begin{problem}{3.1.11}
悬挂在原点的“单摆”，两边用旋轮线

\[
x = a (\varphi - \sin \varphi), y = a (1 - \cos \varphi)
\]

(其中 $\varphi$ 为产生旋轮线的半径为 a 的圆轮转过的角度)做成挡板来对“单摆”的摆长加以限制,如图所示.“单摆”由长度为 4a 的不可伸长的轻绳和摆锤构成.证明:

(1) 摆锤的轨迹是一条旋轮线；

(2) 摆锤的运动是周期运动, 其周期为 $T = 2\pi \sqrt{\frac{4a}{g}}$ , 与摆幅无关.

提示:用自然坐标的运动微分方程,证明 $\sin\theta$ 做简谐振动, $\theta$ 是摆线的直线部分与 y 轴的夹角.
\begin{center}
\begin{tikzpicture}[>=Stealth, scale=1]
  \def\a{0.44}

  % Coordinate axes
  \draw[->, thick] ({-2*pi*\a - 0.30}, 0) -- ({2*pi*\a + 0.35}, 0) node[below=2pt] {$x$};
  \draw[->, thick] (0, 0.42) -- (0, {-4*\a - 0.55}) node[right=2pt] {$y$};
  \node[above left=1pt] at (0, 0) {$O$};

  % Left and right cycloid barriers
  \draw[thick, domain=0:360, samples=72, smooth]
    plot ({-\a*(\x*pi/180 - sin(\x))}, {-\a*(1 - cos(\x))});
  \draw[thick, domain=0:360, samples=72, smooth]
    plot ({\a*(\x*pi/180 - sin(\x))}, {-\a*(1 - cos(\x))});

  % Bob trajectory (involute cycloid)
  \draw[dashed, thick, domain=-180:180, samples=72, smooth]
    plot ({\a*(\x*pi/180 + sin(\x))}, {-\a*(3 + cos(\x))});

  % Length 4a annotation
  \draw[thin] (-0.42, {-4*\a}) -- (0.12, {-4*\a});
  \draw[<->, thin] (-0.24, 0) -- (-0.24, {-4*\a}) node[pos=0.58, left=2pt, fill=white, inner sep=1pt] {$4a$};
\end{tikzpicture}

\end{center}
\end{problem}
\clearpage


\begin{problem}{3.1.12}
若放置在光滑的水平面上的弹簧振动系统,弹簧的质量不可忽略,证明系统的振动周期为

\[
T = 2 \pi \sqrt {\frac {M + m / 3}{k}}
\]

其中 m 为弹簧的质量, k 为弹簧的劲度系数, M 为系于弹簧上的物体的质量.

提示: 可将弹簧等效于一个弹簧振子: 质量可忽略的弹簧的原长和劲度系数均与质量不可忽略的弹簧一样, 连在弹簧上的物体质量为 $m'$ , 先证明 $m' = \frac{1}{3}m$ .
\end{problem}
\clearpage


\begin{problem}{3.1.13}
一个质量为 $m$ 的小物体放在光滑的水平面上被两个原长为 $l_0$ 、劲度系数为 $k$ 的弹簧约束，在图3.11所示的平衡位置附近做小振动。

(1)求小振动的角频率；

(2) 如果弹簧原长变为 $2l_{0}$ ，运动将会发生什么定性变化？
\begin{center}
\begin{tikzpicture}[>=Stealth, scale=1]
  % Hatched walls and floor
  \fill[pattern=north east lines] (-0.22, -0.22) rectangle (0, 1.35);
  \fill[pattern=north east lines] (0, -0.22) rectangle (5.2, 0);
  \fill[pattern=north east lines] (5.2, -0.22) rectangle (5.42, 1.35);
  \draw[thick] (0, 1.35) -- (0, 0) -- (5.2, 0) -- (5.2, 1.35);

  % Mass block
  \draw[thick, fill=white] (2.25, 0) rectangle (2.95, 0.85);
  \node[above=2pt] at (2.60, 0.85) {$m$};

  % Left and right springs
  \draw[thick, decorate, decoration={zigzag, segment length=5.5, amplitude=3.2, pre length=14, post length=14}] (0, 0.425) -- (2.25, 0.425);
  \draw[thick, decorate, decoration={zigzag, segment length=5.5, amplitude=3.2, pre length=14, post length=14}] (2.95, 0.425) -- (5.2, 0.425);
  \node at (1.125, 0.72) {$k$};
  \node at (4.075, 0.72) {$k$};

  % Dimension labels 2l_0
  \draw[thin] (0, -0.28) -- (0, -0.96);
  \draw[thin] (2.60, -0.28) -- (2.60, -0.96);
  \draw[thin] (5.2, -0.28) -- (5.2, -0.96);
  \draw[<->, thin] (0, -0.78) -- (2.60, -0.78) node[midway, above=2pt] {$2l_0$};
  \draw[<->, thin] (2.60, -0.78) -- (5.2, -0.78) node[midway, above=2pt] {$2l_0$};
\end{tikzpicture}

\end{center}
\end{problem}
\clearpage


\begin{problem}{3.1.14}
一个机器内某零件的振动规律为

\[
x = 0. 4 \sin \omega t + 0. 3 \cos \omega t
\]

x 的单位为 cm, $\omega = 20s^{-1}$ . 求这个振动的振幅、最大速度及最大加速度.

\[
x = A \cos (\omega t + \alpha) = A \cos \omega t \cos \alpha - A \sin \omega t \sin \alpha
\]

与 $x = 0.4\sin \omega t + 0.3\cos \omega t$ 比较，可得

\[
A \mathrm{cos} \alpha = 0. 3, \quad A \mathrm{sin} \alpha = - 0. 4
\]

解得

\[
A = \sqrt {(0 . 3) ^ {2} + (- 0 . 4) ^ {2}} = 0. 5 (\mathrm{cm})
\]

\[
v _ {\max} = A w = 0. 5 \times 2 0 = 1 0 (\mathrm{cm} \cdot \mathrm{s} ^ {- 1})
\]

\[
a _ {\max} = A w ^ {2} = 0. 5 \times (2 0) ^ {2} = 2 0 0 (\mathrm{cm} \cdot \mathrm{s} ^ {- 2})
\]
\end{problem}
\clearpage


\begin{problem}{3.1.15}
判断下列振动是否周期运动,若是周期运动,求出它的周期.

期 $T = \sqrt{2}\pi .$
\end{problem}
\clearpage


\begin{problem}{3.1.16}
两个在同一直线上做频率相同的简谐振动的合振动的振幅为 $10\mathrm{cm}$ , 合振动的相位超前第一个振动 $\frac{\pi}{6}$ , 第一个振动的振幅 $A_{1} = 8.0\mathrm{cm}$ . 求第二个振动的振幅 $A_{2}$ 及它与第一个振动的相位差.
\end{problem}
\clearpage


\begin{problem}{3.1.17}
一弹簧振子的劲度系数 $k_{1} = 9.8\mathrm{N}/$ $\mathbf{m}$ ，质量 $m = 9.8\times 10^{-2}\mathrm{kg}$ ，它的影子水平地投射在一屏上，该屏的质量 $M = 0.98\mathrm{kg}$ ，通过劲度系数 $k_{2}$ $= 98\mathrm{N / m}$ 的弹簧挂起来，如图3.13所示．开始时，把它们都从平衡位置拉下 $0.10\mathrm{m}$ ，先释放弹簧振子，如果要使影子在屏上振动的振幅为 $0.050\mathrm{m}$ ，问要过多长时间释放屏？并写出影子在屏上的振动方程.
\begin{center}
\begin{tikzpicture}[>=Stealth, scale=1]
  % Hatched ceiling
  \fill[pattern=north east lines] (-0.2, 3.2) rectangle (4.5, 3.42);
  \draw[thick] (-0.2, 3.2) -- (4.5, 3.2);

  % Left spring and mass m
  \draw[thick, decorate, decoration={zigzag, segment length=5.5, amplitude=3.2, pre length=14, post length=8}] (0.9, 3.2) -- (0.9, 1.51);
  \node[left=5pt] at (0.9, 2.35) {$k_1$};
  \draw[thick, fill=white] (0.9, 1.15) circle (0.36);
  \node[left=3pt] at (0.54, 1.15) {$m$};

  % Right spring and screen M
  \draw[thick, decorate, decoration={zigzag, segment length=5.5, amplitude=3.2, pre length=6, post length=6}] (3.25, 3.2) -- (3.25, 2.00);
  \node[right=5pt] at (3.25, 2.65) {$k_2$};
  \draw[thick, fill=white] (2.45, 2.25) -- (4.05, 1.75) -- (4.05, 0.05) -- (2.45, 0.55) -- cycle;
  \node at (3.65, 0.55) {$M$};

  % Horizontal projection onto screen
  \draw[dashed, thick] (0.9, 1.51) -- (3.25, 1.51);
  \draw[dashed, thick] (0.9, 0.79) -- (3.25, 0.79);
  \draw[dashed, thick] (3.25, 1.15) ellipse (0.14 and 0.36);
\end{tikzpicture}

\end{center}
\end{problem}
\clearpage


\begin{problem}{3.1.18}
用旋转矢量法证明，在同一直线上的 N 个简谐振动，它们有相同的振幅 $A_{0}$ 、相同的角频率 $\omega_{0}$ 、相位依次超前 $\delta$ ，即

\[
x _ {i} = A _ {0} \cos [ \omega t + (i - 1) \delta ] \quad (i = 1, 2, \dots , N)
\]

叠加后的合振动

\[
x = \sum_ {i = 1} ^ {N} x _ {i} = A _ {0} \frac {\sin \left(\frac {1}{2} N \delta\right)}{\sin \left(\frac {1}{2} \delta\right)} \cos \left[ \omega t + \frac {1}{2} (N - 1) \delta \right]
\]

提示:考虑 N 个振幅矢量都是一个圆上的弦,每个振幅矢量是正多边形的一条边.
\end{problem}
\clearpage


\begin{problem}{3.1.19}
有两支标有 $256\mathrm{Hz}$ 的音叉, 一支是标准的, 另一支是待校正的. 同时敲击两支音叉, 在半分钟内能听到 45 拍, 若在待校正的音叉上适当涂蜡, 从少逐渐增多, 在半分钟内听到的拍数先减少到听不见拍声, 而后又增多, 问待校正音叉原来的频率多大?
\end{problem}
\clearpage


\begin{problem}{3.1.20}
画出在互相垂直的两个方向振动的合成运动的轨道,并在图上标明运动方向.

\[
x = A \sin \omega t, y = B \sin \left(\omega t - \frac {\pi}{3}\right); \tag {1}
\]

(2) $x = A\sin \omega t, y = B\sin \left(2\omega t + \frac{5}{4}\pi\right)$ .
\end{problem}
\clearpage


\begin{problem}{3.1.21}
图 3.18 中的各图是在相互垂直的两个方向振动合成的利萨如图形: 已知水平方向的振动的角频率为 $\omega$ , 求竖直方向的振动的角频率.

(a)

(b)

(c)
\begin{center}
\begin{tikzpicture}[>=Stealth, scale=1]
  \draw[thick] (-1.55, -2.05) rectangle (1.55, 2.05);
  \draw[thick, domain=0:360, samples=240, smooth]
    plot ({1.55*sin(3*\x)}, {2.05*sin(\x + 15)});
\end{tikzpicture}

\end{center}
\begin{center}
\begin{tikzpicture}[>=Stealth, scale=1]
  \draw[thick] (-1.55, -2.05) rectangle (1.55, 2.05);
  \draw[thick, domain=0:180, samples=200, smooth]
    plot ({-1.55*cos(4*\x)}, {2.05*cos(3*\x)});
\end{tikzpicture}

\end{center}
\begin{center}
\begin{tikzpicture}[>=Stealth, scale=1]
  \draw[thick] (-1.55, -2.05) rectangle (1.55, 2.05);
  \draw[thick, domain=0:360, samples=240, smooth]
    plot ({1.55*sin(3*\x)}, {-2.05*cos(2*\x - 20)});
\end{tikzpicture}

\end{center}
\end{problem}
\clearpage


\begin{problem}{3.1.22}
在图 3.19 所示的示波器荧光屏上显示的利萨如图形中, 已知水平偏转讯号的频率为 $100\mathrm{Hz}$ , 若运行方向如图中箭头所示, 且 $t = 0$ 时位于 $S$ 点. 试写出水平偏转讯号和竖直偏转讯号的表示式, 并求出 $P, Q$ 两点的 $x$ 坐标.
\begin{center}
\begin{tikzpicture}[>=Stealth, line cap=round, line join=round]
  \def\A{2.0}
  \def\B{2.0}

  % Bounding rectangle
  \draw[thick] (-\A, -\B) rectangle (\A, \B);

  % Coordinate axes
  \draw[->, thick] (-\A, 0) -- (\A + 0.55, 0) node[below] {$x$};
  \draw[->, thick] (0, -\B) -- (0, \B + 0.55) node[left] {$y$};

  % Lissajous curve: x(t) = A cos(2t), y(t) = -B cos(3t) for t in [0, pi]
  \draw[thick,
    postaction={decorate},
    decoration={markings, mark=at position 0.21 with {\arrow{Stealth}}}
  ] plot[domain=0:3.14159265, samples=80, smooth]
    ({\A*cos(deg(2*\x))}, {-\B*cos(deg(3*\x))});

  % Key point labels
  \node[below right] at (\A, -\B) {$S$};
  \node[below right] at (0.5*\A, 0) {$P$};
  \node[below] at (-0.5*\A, \B - 0.08) {$Q$};

  % Dimension 2B (left)
  \draw (-\A - 0.15, \B) -- (-\A - 0.55, \B);
  \draw (-\A - 0.15, -\B) -- (-\A - 0.55, -\B);
  \draw[<->] (-\A - 0.35, -\B) -- (-\A - 0.35, \B)
    node[midway, left] {$2B$};

  % Dimension 2A (bottom)
  \draw (-\A, -\B - 0.15) -- (-\A, -\B - 0.55);
  \draw (\A, -\B - 0.15) -- (\A, -\B - 0.55);
  \draw[<->] (-\A, -\B - 0.35) -- (\A, -\B - 0.35)
    node[midway, below] {$2A$};
\end{tikzpicture}

\end{center}
\end{problem}
\clearpage


\begin{problem}{3.1.23}
上题中, 若 t=0 时位于 P 点, 写出水平、竖直两方向偏转讯号的表示式.
\end{problem}
\clearpage


\begin{problem}{3.1.24}
对图 3.20 所示的各周期运动进行谐波分析.

(a)

(b)

(c)
\begin{center}
\begin{tikzpicture}[>=Stealth, line cap=round, line join=round]
  \def\T{1.3}
  \def\A{1.4}

  % Axes
  \draw[->, thick] (-0.35, 0) -- (3.45*\T, 0) node[below] {$t$};
  \draw[->, thick] (0, -0.35) -- (0, \A + 0.55) node[left] {$x$};
  \node[below left] at (0, 0) {$O$};
  \node[left] at (0, \A) {$A$};

  % Sawtooth periods
  \foreach \k in {0, 1, 2} {
    \draw[thick] (\k*\T, 0) -- ({(\k+1)*\T}, \A);
    \draw[thick, dashed] ({(\k+1)*\T}, \A) -- ({(\k+1)*\T}, 0);
  }
  % Continuation after 3rd period
  \draw[thick] (3*\T, 0) -- (3.22*\T, 0.22*\A);

  % Period labels
  \node[below] at (\T, -0.05) {$\dfrac{2\pi}{\omega}$};
  \node[below] at (2*\T, -0.05) {$\dfrac{4\pi}{\omega}$};
  \node[below] at (3*\T, -0.05) {$\dfrac{6\pi}{\omega}$};
\end{tikzpicture}

\end{center}
\begin{center}
\begin{tikzpicture}[>=Stealth, line cap=round, line join=round]
  \def\qT{0.75} % quarter period pi/(2 omega)
  \def\A{1.05}

  % Axes
  \draw[->, thick] (-0.4, 0) -- (6.1*\qT, 0) node[below] {$t$};
  \draw[->, thick] (0, -\A - 0.35) -- (0, \A + 0.55) node[left] {$x$};
  \node[below left] at (0, 0) {$O$};
  \node[left] at (0, \A) {$A$};
  \node[left] at (0, -\A) {$-A$};

  % Square wave
  \draw[thick] (0, \A) -- (\qT, \A);
  \draw[thick, dashed] (\qT, \A) -- (\qT, -\A);
  \draw[thick] (\qT, -\A) -- (3*\qT, -\A);
  \draw[thick, dashed] (3*\qT, -\A) -- (3*\qT, \A);
  \draw[thick] (3*\qT, \A) -- (5*\qT, \A);
  \draw[thick, dashed] (5*\qT, \A) -- (5*\qT, -\A);
  \draw[thick] (5*\qT, -\A) -- (5.55*\qT, -\A);

  % Ticks and labels on t-axis
  \draw (2*\qT, 0) -- (2*\qT, 0.08);
  \node[below] at (2*\qT, -0.05) {$\dfrac{\pi}{\omega}$};
  \draw (4*\qT, 0) -- (4*\qT, 0.08);
  \node[below] at (4*\qT, -0.05) {$\dfrac{2\pi}{\omega}$};
\end{tikzpicture}

\end{center}
\begin{center}
\begin{tikzpicture}[>=Stealth, line cap=round, line join=round]
  \def\hT{0.82} % half period pi/omega
  \def\A{1.35}

  % Axes
  \draw[->, thick] (-0.35, 0) -- (5.7*\hT, 0) node[below] {$t$};
  \draw[->, thick] (0, -0.35) -- (0, \A + 0.55) node[left] {$x$};
  \node[below left] at (0, 0) {$O$};
  \node[left] at (0, \A) {$A$};

  % Triangle wave
  \draw[thick] (0, \A) -- (\hT, 0) -- (2*\hT, \A) -- (3*\hT, 0)
    -- (4*\hT, \A) -- (5*\hT, 0) -- (5.25*\hT, 0.25*\A);

  % Ticks and labels on t-axis
  \node[below] at (\hT, -0.05) {$\dfrac{\pi}{\omega}$};
  \draw (2*\hT, 0) -- (2*\hT, 0.07);
  \node[below] at (2*\hT, -0.05) {$\dfrac{2\pi}{\omega}$};
  \node[below] at (3*\hT, -0.05) {$\dfrac{3\pi}{\omega}$};
\end{tikzpicture}

\end{center}
\end{problem}
\clearpage


\section{阻尼振动和受迫振动}

\begin{problem}{3.2.1}
一质量为 $m$ 的质点在恢复力 $-kx$ 和阻力 $-r\dot{x}(k$ 和 $r$ 均为正的常量)的作用下运动, 其中 $x$ 是质点偏离平衡位置的位移. 对于固定的 $k$ 和一般的初始条件, $r = r_0$ 为何值时回到平衡位置所需的时间最短. 如果 $r > r_0$ , 是否可能对于特定的初始条件, 比 $r = r_0$ 能更快回到平衡位置? 如果 $r < r_0$ 呢?
\end{problem}
\clearpage


\begin{problem}{3.2.2}
如图 3.21 所示, 质量为 m=2000kg 的重物以 v=3cm/s 的速度匀速运动, 与质量可以忽略不计的弹簧及阻尼器相碰撞后一起做自由阻尼振动. 已知弹簧的劲度系数 k=48020N/m, 阻尼

器的阻尼系数 $\gamma = 1960\mathrm{N}\mathrm{s / m}$ ，问重物在碰撞后多少时间达到最大位移？最大位移多大？解沿 $v$ 的方向取 $x$ 轴，与阻尼器刚碰撞处取为 $x = 0$ ，此时 $t = 0$ 

\[
m \dot {x} = - k x - \gamma \dot {x}
\]

令 $\beta = \frac{\gamma}{2m},\omega_0^2 = \frac{k}{m}$ ，上式可改写为

\[
\dot {x} + 2 \beta \dot {x} + \omega_ {0} ^ {2} x = 0
\]

相应的特征方程为

\[
\begin{array}{r l} & r ^ {2} + 2 \beta r + \omega_ {0} ^ {2} = 0 \\ & r = - \beta \pm \sqrt {\beta^ {2} - \omega_ {0} ^ {2}} \\ & \omega_ {0} = \sqrt {\frac {k}{m}} = \sqrt {\frac {4 8 0 2 0}{2 0 0 0}} = 4. 9 \mathrm{s} ^ {- 1} \\ & \beta = \frac {\gamma}{2 m} = \frac {1 9 6 0}{2 \times 2 0 0 0} = 0. 4 9 \mathrm{s} ^ {- 1} \end{array}
\]

$\beta<\omega_{0}$ ，为弱阻尼情况，通解为

\[
x = A e ^ {- \beta t} \cos (\sqrt {\omega_ {0} ^ {2} - \beta^ {2}} t + \alpha)
\]

初始条件: t=0 时, x=0, $\dot{x}=v$ .

解得

\[
\begin{array}{r l} \dot {x} & = - \beta A \mathrm{e} ^ {- \beta t} \cos (\sqrt {\omega_ {0} ^ {2} - \beta^ {2}} t + \alpha) \\ & - A \sqrt {\omega_ {0} ^ {2} - \beta^ {2}} \mathrm{e} ^ {- \beta t} \sin (\sqrt {\omega_ {0} ^ {2} - \beta^ {2}} t + \alpha) \\ & A \cos \alpha = 0 \\ & - \beta A \cos \alpha - A \sqrt {\omega_ {0} ^ {2} - \beta^ {2}} \sin \alpha = v \\ & \alpha = - \frac {\pi}{2}, A = \frac {v}{\sqrt {\omega_ {0} ^ {2} - \beta^ {2}}} \end{array}
\]

\[
\begin{array}{c} x = \frac {v}{\sqrt {\omega_ {0} ^ {2} - \beta^ {2}}} \mathrm{e} ^ {- \beta t} \sin (\sqrt {\omega_ {0} ^ {2} - \beta^ {2}} t) \\ \dot {x} = - \frac {\beta v}{\sqrt {\omega_ {0} ^ {2} - \beta^ {2}}} \mathrm{e} ^ {- \beta t} \sin (\sqrt {\omega_ {0} ^ {2} - \beta^ {2}} t) + v \mathrm{e} ^ {- \beta t} \cos (\sqrt {\omega_ {0} ^ {2} - \beta^ {2}} t) \end{array}
\]

重物达到最大位移时， $\dot{x} = 0$ ，此时 $t = t_1$ 

\[
\begin{array}{r l} & {- \frac {\beta v}{\sqrt {\omega_ {0} ^ {2} - \beta^ {2}}} \sin (\sqrt {\omega_ {0} ^ {2} - \beta^ {2}} t _ {1}) + v \cos (\sqrt {\omega_ {0} ^ {2} - \beta^ {2}} t _ {1}) = 0} \\ & {\qquad t _ {1} = \frac {1}{\sqrt {\omega_ {0} ^ {2} - \beta^ {2}}} \arctan \left(\frac {\sqrt {\omega_ {0} ^ {2} - \beta^ {2}}}{\beta}\right) = 0. 3 0 \mathrm{s}} \end{array}
\]

最大位移为

\[
x \left(t _ {1}\right) = \frac {v}{\sqrt {\omega_ {0} ^ {2} - \beta^ {2}}} \mathrm{e} ^ {- \beta t _ {1}} \sin \left(\sqrt {\omega_ {0} ^ {2} - \beta^ {2}} t _ {1}\right) = 5. 2 8 \times 1 0 ^ {- 3} \mathrm{m}
\]
\begin{center}
\begin{tikzpicture}[>=Stealth, line cap=round, line join=round]
  % Left wall with hatching
  \fill[pattern=north east lines] (-0.25, 0.35) rectangle (0, 2.25);
  \draw[thick] (0, 0.35) -- (0, 2.25);

  % Horizontal ground with hatching
  \fill[pattern=north east lines] (0, -0.22) rectangle (5.6, 0);
  \draw[thick] (0, 0) -- (5.6, 0);

  % Dashpot (top)
  \draw[thick] (0, 1.65) -- (1.2, 1.65);
  \draw[thick] (1.2, 1.42) -- (1.2, 1.88);
  \draw[thick] (0.95, 1.98) -- (1.55, 1.98) -- (1.55, 1.32) -- (0.95, 1.32);
  \draw[thick] (1.55, 1.65) -- (1.92, 1.65);
  \node[above] at (1.25, 1.98) {$\gamma$};

  % Spring (bottom)
  \draw[thick] (0, 0.8) -- (0.5, 0.8);
  \draw[thick, decorate, decoration={zigzag, segment length=7pt, amplitude=3.5pt}]
    (0.5, 0.8) -- (1.5, 0.8);
  \draw[thick] (1.5, 0.8) -- (1.92, 0.8);
  \node[below] at (1.0, 0.62) {$k$};

  % Vertical bumper plate
  \draw[thick] (1.92, 0.35) rectangle (2.12, 2.05);

  % Velocity arrow
  \draw[->, thick] (3.55, 1.25) -- (2.8, 1.25) node[midway, above] {$v$};

  % Mass cart with wheels
  \draw[thick] (3.75, 0.42) rectangle (5.3, 1.65);
  \node at (4.525, 1.035) {$m$};
  \draw[thick] (4.1, 0.21) circle (0.21);
  \draw[thick] (4.95, 0.21) circle (0.21);
\end{tikzpicture}

\end{center}
\end{problem}
\clearpage


\begin{problem}{3.2.3}
一个有阻尼的弹簧振子系统, $m = 100\mathrm{kg}, k = 10000\mathrm{N / m}$ , 处于临界阻尼情况,由 $t = 0$ 时, $x_0 = 2.5\mathrm{cm}, \dot{x}_0 = -30\mathrm{cm / s}$ 开始运动. 问质量块将于几秒后到达平衡位置? 达到平衡位置后最远能移动多大距离?
\end{problem}
\clearpage


\begin{problem}{3.2.4}
如图 3.22 所示, 质量为 m 的薄板挂在弹簧的下端, 若空气的阻力可以不计,在空气中的振动周期为 $T_{1}$ ，在某液体中的振动周期为 $T_{2}$ ，液体的阻力可表示为 $-2\eta A v$ ，其中 A 为薄板的面积，v 为速度，

$\eta$ 为液体的黏度，试证 $\eta = \frac{2\pi m}{AT_1T_2}\cdot \sqrt{T_2^2 - T_1^2}.$
\begin{center}
\begin{tikzpicture}[>=Stealth, line cap=round, line join=round]
  % Ceiling
  \fill[pattern=north east lines] (-1.6, 3.6) rectangle (1.6, 3.85);
  \draw[thick] (-1.6, 3.6) -- (1.6, 3.6);

  % Spring and suspension rod
  \draw[thick] (0, 3.6) -- (0, 3.25);
  \draw[thick, decorate, decoration={zigzag, segment length=7pt, amplitude=3.5pt}]
    (0, 3.25) -- (0, 2.3);
  \node[right] at (0.12, 2.78) {$k$};
  \draw[thick] (0, 2.3) -- (0, 1.75);

  % Liquid surface
  \draw[thick] (-1.9, 2.2) -- (1.9, 2.2);

  % Liquid dashes (around plate)
  \foreach \x/\y in {-1.6/1.9, -1.0/1.85, 0.7/1.9, 1.4/1.85,
                     -1.5/1.45, -0.95/1.4, 1.05/1.45, 1.6/1.4,
                     -1.65/1.0, -1.05/0.95, 1.05/1.0, 1.6/0.95,
                     -1.5/0.55, -0.95/0.5, 1.1/0.55, 1.55/0.5,
                     -1.55/0.1, -0.85/0.05, 1.05/0.1, 1.55/0.05} {
    \draw[thick] (\x, \y) -- (\x + 0.22, \y);
  }

  % 3D thin rectangular plate (submerged)
  % Front face: (-0.5, 1.85), (0.55, 1.45), (0.55, 0.05), (-0.5, 0.45)
  % Back-top offset: (0.18, 0.1)
  \fill[white] (-0.5, 1.85) -- (0.55, 1.45) -- (0.55, 0.05) -- (-0.5, 0.45) -- cycle;
  \draw[thick] (-0.5, 1.85) -- (0.55, 1.45) -- (0.55, 0.05) -- (-0.5, 0.45) -- cycle;
  % Side and top thickness edges
  \draw[thick] (-0.5, 1.85) -- (-0.32, 1.95) -- (0.73, 1.55) -- (0.73, 0.15) -- (0.55, 0.05);
  \draw[thick] (0.55, 1.45) -- (0.73, 1.55);

  % Suspension rod visible down to attachment point on top face
  \draw[thick] (0, 2.2) -- (0, 1.71);

  % Labels on plate
  \node at (0.02, 1.30) {$m$};
  \node at (0.02, 0.78) {$A$};
\end{tikzpicture}

\end{center}
\end{problem}
\clearpage


\begin{problem}{3.2.5}
在图 3.23 所示系统中, 刚杆质量可忽略不计, 图中质量 m、弹簧的劲度系数 k、阻尼器阻尼系数 $\gamma$ 以及距离 a、b、l 均为已知. 写出运动微分方程. 由此求出做弱阻尼振动的角频率 $\omega_{d}$ 以及临界阻尼时 $\gamma = \gamma_{cr}$ 与其余各量间必须满足的关系.
\begin{center}
\begin{tikzpicture}[>=Stealth, line cap=round, line join=round]
  % Wall and pivot hinge at (0, 0)
  \fill[pattern=north east lines] (-0.42, -0.45) rectangle (-0.22, 0.45);
  \draw[thick] (-0.22, -0.45) -- (-0.22, 0.45);
  \draw[thick] (-0.22, 0.24) -- (0.0, 0.24) arc (90:-90:0.24) -- (-0.22, -0.24);
  \draw[thick] (0, 0) circle (0.06);

  % Rigid horizontal rod (from x=0.06 to x=4.68) and mass bob at x=4.95
  \draw[thick] (0.06, 0.07) -- (4.72, 0.07);
  \draw[thick] (0.06, -0.07) -- (4.72, -0.07);
  \draw[thick] (4.72, 0.07) arc (166:-166:0.25) -- cycle;
  \node[above right] at (4.95, 0.18) {$m$};

  % Dashpot at x = 1.9 (above rod)
  \fill[pattern=north east lines] (1.6, 1.75) rectangle (2.2, 1.95);
  \draw[thick] (1.6, 1.75) -- (2.2, 1.75);
  \draw[thick] (1.9, 1.75) -- (1.9, 1.15);
  \draw[thick] (1.72, 1.15) -- (2.08, 1.15);
  \draw[thick] (1.65, 1.35) -- (1.65, 0.85) -- (2.15, 0.85) -- (2.15, 1.35);
  \draw[thick] (1.9, 0.85) -- (1.9, 0.07);
  \node[right] at (2.18, 1.1) {$\gamma$};

  % Spring at x = 3.2 (below rod)
  \draw[thick] (3.2, -0.07) -- (3.2, -0.45);
  \draw[thick, decorate, decoration={zigzag, segment length=7pt, amplitude=3.5pt}]
    (3.2, -0.45) -- (3.2, -1.35);
  \draw[thick] (3.2, -1.35) -- (3.2, -1.75);
  \node[right] at (3.32, -0.9) {$k$};
  \fill[pattern=north east lines] (2.9, -1.95) rectangle (3.5, -1.75);
  \draw[thick] (2.9, -1.75) -- (3.5, -1.75);

  % Dimension extension lines
  \draw (0, -0.35) -- (0, -3.05);
  \draw (1.9, -0.2) -- (1.9, -1.0);
  \draw (3.2, -2.0) -- (3.2, -2.4);
  \draw (4.95, -0.35) -- (4.95, -3.05);

  % Dimension arrows a, b, l
  \draw[<->] (0, -0.75) -- (1.9, -0.75) node[midway, above] {$a$};
  \draw[<->] (0, -2.25) -- (3.2, -2.25) node[midway, above] {$b$};
  \draw[<->] (0, -2.88) -- (4.95, -2.88) node[midway, above] {$l$};
\end{tikzpicture}

\end{center}
\end{problem}
\clearpage


\begin{problem}{3.2.6}
质量为 0.2kg 的物体, 悬于劲度系数为 80N/m 的弹簧下, 此物体还受到阻力 $-\gamma v$ 的作用, v 是物体的速度. 若阻尼振动的频率为无阻尼振动频率的 $\frac{\sqrt{3}}{2}$ 倍, 阻尼系数 $\gamma$ 多大? 此系统 Q 值多大? 振动经 10 个整周后, 振幅减弱到原振幅的几分之几?
\end{problem}
\clearpage


\begin{problem}{3.2.7}
演奏钢琴上频率为 256Hz 的音键时, 1s 后振动能量减至初始值的一半. 问系统的 Q 值多大? 若高八度的音调(512Hz)也在 1s 后振动能量减至初始值的一半, 其 Q 值多大?
\end{problem}
\clearpage


\begin{problem}{3.2.8}
根据经典电磁理论,一个加速运动的电子在单位时间内辐射的能量为 $ke^{2}a^{2}/c^{3}$ , 其中 $k=6\times10^{9}Nm^{2}/C^{2}, e=1.60\times10^{-19}C$ (注意 k、e 的单位中的 C 是库仑, 勿与光速 c 混淆) $c=3.0\times10^{8}m/s, a$ 为电子的瞬时加速度.

(1) 若一电子沿一直线以频率 $\nu$ 及振幅 A 进行振荡, 假定每一周的运动可用 $x = A \cos(2\pi \nu t)$ 描述, 问在振荡一周的时间内, 它辐射了多少能量?

(2) 此振子之 Q 值多大?

(3) 在电子的运动的能量降至初始值的一半以前, 振荡已经历了多少个周期?
\end{problem}
\clearpage


\begin{problem}{3.2.9}
有一阻尼振子, 质量 $m = 0.2 \mathrm{~kg}, k = 80 \mathrm{~N} / \mathrm{m}, \gamma = 4 \mathrm{~N} \mathrm{~s} / \mathrm{m}$ , 受到强迫力 $F = F_{0} \cos \Omega t$ 作用, $F_{0} = 2 \mathrm{~N}, \Omega = 30 \mathrm{~s}^{-1}$ , 求:

(1) 稳态运动 $x = B \cos(\Omega t - \phi)$ 中的 B 和 $\phi$ ;

(2) 一周内反抗阻力耗散的机械能；

(3) 输入的平均功率.
\end{problem}
\clearpage


\begin{problem}{3.2.10}
如图 3.25 所示,作用于质量块上的强迫力 $F(t) = F_0\sin \Omega t$ , 弹簧支承端有运动, $x_s = a\cos \Omega t, \Omega \neq \sqrt{\frac{k}{m}}$ , $k$ 为弹簧的劲度系数, $m$ 为质量块的质量, 求稳态振动.
\begin{center}
\begin{tikzpicture}[>=Stealth, line cap=round, line join=round]
  % Left guide walls for sliding support
  \fill[pattern=north east lines] (0, 1.55) rectangle (1.15, 1.78);
  \draw[thick] (0, 1.55) -- (1.15, 1.55);
  \fill[pattern=north east lines] (0, 0.22) rectangle (1.15, 0.45);
  \draw[thick] (0, 0.45) -- (1.15, 0.45);

  % Sliding support block
  \draw[thick] (0.15, 0.52) rectangle (0.85, 1.48);

  % Spring
  \draw[thick] (0.85, 1.0) -- (1.25, 1.0);
  \draw[thick, decorate, decoration={zigzag, segment length=7pt, amplitude=3.5pt}]
    (1.25, 1.0) -- (2.45, 1.0);
  \draw[thick] (2.45, 1.0) -- (2.85, 1.0);
  \node[above] at (1.85, 1.18) {$k$};

  % Floor under cart
  \fill[pattern=north east lines] (2.2, -0.22) rectangle (4.65, 0);
  \draw[thick] (2.2, 0) -- (4.65, 0);

  % Mass cart with wheels
  \draw[thick] (2.85, 0.42) rectangle (4.15, 1.45);
  \node at (3.5, 0.935) {$m$};
  \draw[thick] (3.15, 0.21) circle (0.21);
  \draw[thick] (3.85, 0.21) circle (0.21);

  % Applied force F(t)
  \draw[->, thick] (4.15, 0.95) -- (5.2, 0.95) node[midway, above] {$F(t)$};

  % Coordinate axes below
  \draw[->, thick] (0.15, -0.65) -- (1.45, -0.65) node[below] {$x_s$};
  \node[below] at (0.15, -0.68) {$0$};

  \draw[->, thick] (2.25, -0.65) -- (5.2, -0.65) node[below] {$x$};
  \draw (2.85, -0.65) -- (2.85, -0.57);
  \node[below] at (2.85, -0.68) {$0$};
\end{tikzpicture}

\end{center}
\end{problem}
\clearpage


\begin{problem}{3.2.11}
图 3.28 所示的系统中, 弹簧悬点上下做简谐振动, $x_{s} = \frac{v_{0}}{\omega}\sin \omega t$ , 初始条件为 $x(0) = 0$ , $\dot{x}(0) = v_{0}$ , 弹簧的劲度系数为 $k$ , 质量块的质量为 $m$ , $\omega = \sqrt{\frac{k}{m}}$ . 求质量块的运动学方程.
\begin{center}
\begin{tikzpicture}[>=Stealth, line cap=round, line join=round]
  % Guide walls on left and right of sliding support
  \fill[pattern=north east lines] (-1.35, 3.0) rectangle (-0.45, 3.6);
  \draw[thick] (-1.35, 3.6) -- (-0.45, 3.6) -- (-0.45, 3.0) -- (-1.35, 3.0);

  \fill[pattern=north east lines] (0.45, 3.0) rectangle (1.35, 3.6);
  \draw[thick] (1.35, 3.6) -- (0.45, 3.6) -- (0.45, 3.0) -- (1.35, 3.0);

  % Sliding support block
  \draw[thick] (-0.38, 2.88) rectangle (0.38, 3.72);

  % Spring
  \draw[thick] (0, 2.88) -- (0, 2.38);
  \draw[thick, decorate, decoration={zigzag, segment length=7pt, amplitude=3.5pt}]
    (0, 2.38) -- (0, 1.38);
  \node[left] at (-0.15, 1.88) {$k$};
  \draw[thick] (0, 1.38) -- (0, 0.88);

  % Mass block
  \draw[thick] (-0.65, 0.25) rectangle (0.65, 0.88);
  \node[left] at (-0.72, 0.56) {$m$};

  % Coordinate axes on the right
  \draw[->, thick] (1.75, 3.72) -- (1.75, 2.35) node[right] {$x_s$};
  \node[right] at (1.78, 3.62) {$0$};

  \draw[->, thick] (0.95, 2.3) -- (0.95, -0.25) node[right] {$x$};
  \node[right] at (0.98, 2.2) {$0$};
\end{tikzpicture}

\end{center}
\end{problem}
\clearpage


\begin{problem}{3.2.12}
如图 3.29 所示, 单摆悬点沿水平方向做简谐振动, $x_{s}=a\sin\omega t$ , 已知摆长为 l, 摆锤质量为 m, 求微幅稳态振动中 $\theta$ 随时间的变化规律.
\begin{center}
\begin{tikzpicture}[>=Stealth, line cap=round, line join=round]
  % Horizontal sliding rod
  \draw[thick] (-2.6, 0) -- (2.6, 0);

  % Left bearing guide
  \draw[thick] (-2.2, 0.5) -- (-2.2, 0.1) -- (-1.35, 0.1) -- (-1.35, 0.5);
  \draw[thick] (-2.2, -0.5) -- (-2.2, -0.1) -- (-1.35, -0.1) -- (-1.35, -0.5);

  % Right bearing guide
  \draw[thick] (1.35, 0.5) -- (1.35, 0.1) -- (2.2, 0.1) -- (2.2, 0.5);
  \draw[thick] (1.35, -0.5) -- (1.35, -0.1) -- (2.2, -0.1) -- (2.2, -0.5);

  % Pivot coordinate x_s
  \draw[->, thick] (0, 0.35) -- (0.95, 0.35) node[above] {$x_s$};
  \node[above left] at (0.05, 0.35) {$0$};

  % Vertical dashed reference line
  \coordinate (O) at (0, 0);
  \coordinate (V) at (0, -2.9);
  \draw[thick, dashed] (O) -- (V);

  % Pendulum rod and bob (angle 24 deg from vertical)
  \def\L{2.55}
  \def\ang{24}
  \coordinate (M) at ({\L*sin(\ang)}, {-\L*cos(\ang)});
  \draw[thick] (O) -- node[midway, above right] {$l$} (M);
  \filldraw[fill=white, draw=black, thick] (M) circle (0.18);
  \node[right] at ($(M) + (0.2, 0)$) {$m$};

  % Angle theta
  \pic[draw, angle radius=0.85cm, angle eccentricity=1.35, "$\theta$"] {angle=V--O--M};
\end{tikzpicture}

\end{center}
\end{problem}
\clearpage


\begin{problem}{3.2.13}
图 3.30 所示系统中, 刚性杆 $AB$ 的质量可忽略不计, $B$ 端作用的强迫力 $F(t) = F_0 \sin \omega t$ , 图中 $m, k, \gamma, l$ 均为已知量。写出系统的运动微分方程, 并求下列情况下质量 $m$ 做上下稳态振动的振幅:

(1)系统发生共振；

(2) $\omega$ 等于固有频率 $\omega_0$ 的一半.
\begin{center}
\begin{tikzpicture}[>=Stealth, line width=0.7pt]
  % Left fixed wall and hinge A
  \fill[pattern=north east lines] (-0.65,-0.55) rectangle (-0.45,0.55);
  \draw[line width=0.9pt] (-0.45,-0.55) -- (-0.45,0.55);
  \draw (-0.45,0.22) -- (0,0.08);
  \draw (-0.45,-0.22) -- (0,-0.08);

  % Rigid massless bar AB
  \draw (0,0.08) -- (4.8,0.08) -- (4.8,-0.08) -- (0,-0.08) -- cycle;

  % Hinge circle A and mass m
  \draw[fill=white] (0,0) circle (0.11);
  \node[above=5pt] at (0,0.11) {$A$};
  \draw[fill=white] (1.6,0) circle (0.24);
  \node[above=2pt] at (1.6,0.24) {$m$};

  % Dashpot at x = 3.2
  \draw (3.2,-0.08) -- (3.2,-0.98);
  \draw (3.06,-0.98) -- (3.34,-0.98);
  \draw (2.98,-0.80) -- (2.98,-1.25) -- (3.42,-1.25) -- (3.42,-0.80);
  \draw (3.2,-1.25) -- (3.2,-1.75);
  \node[right=2pt] at (3.42,-1.02) {$\gamma$};
  \fill[pattern=north east lines] (2.8,-1.92) rectangle (3.6,-1.75);
  \draw[line width=0.9pt] (2.8,-1.75) -- (3.6,-1.75);

  % Force F(t) and label B at x = 4.8
  \draw[->] (4.8,0.08) -- (4.8,1.05) node[left=2pt, pos=0.85] {$F(t)$};
  \node[right=3pt] at (4.8,0) {$B$};

  % Spring k at x = 4.8
  \draw (4.8,-0.08) -- (4.8,-0.45);
  \draw[decorate, decoration={coil, aspect=0.4, segment length=4.5pt, amplitude=3.5pt}] (4.8,-0.45) -- (4.8,-1.40);
  \draw (4.8,-1.40) -- (4.8,-1.75);
  \node[right=4pt] at (4.8,-0.92) {$k$};
  \fill[pattern=north east lines] (4.4,-1.92) rectangle (5.2,-1.75);
  \draw[line width=0.9pt] (4.4,-1.75) -- (5.2,-1.75);

  % Bottom dimension lines (l, l, l)
  \foreach \x in {0, 1.6, 3.2, 4.8} {
    \draw[thin] (\x,-2.22) -- (\x,-2.65);
  }
  \draw[<->] (0,-2.48) -- (1.6,-2.48) node[midway, above=2pt] {$l$};
  \draw[<->] (1.6,-2.48) -- (3.2,-2.48) node[midway, above=2pt] {$l$};
  \draw[<->] (3.2,-2.48) -- (4.8,-2.48) node[midway, above=2pt] {$l$};
\end{tikzpicture}

\end{center}
\end{problem}
\clearpage


\begin{problem}{3.2.14}
图 3.31 中所有刚性杆质量均可忽略不计, m、k、 $\gamma$ 和 l 均为已知, 弹簧支承端有运动, $x_{s}=a\sin\omega t$ . 求系统的固有频率 $\omega_{0}$ 、阻尼比 $\zeta\left(=\frac{\beta}{\omega_{0}}\right)$ 及微幅稳态振动的振幅.
\begin{center}
\begin{tikzpicture}[>=Stealth, line width=0.7pt]
  % Top ceiling and hinge
  \fill[pattern=north east lines] (2.25,2.22) rectangle (3.15,2.39);
  \draw[line width=0.9pt] (2.25,2.22) -- (3.15,2.22);
  \draw (2.52,2.22) -- (2.63,1.80);
  \draw (2.88,2.22) -- (2.77,1.80);

  % L-shaped rigid massless bar
  \draw (2.62,1.80) -- (2.62,0.40) -- (2.78,0.40) -- (2.78,1.72) -- (5.40,1.72) -- (5.40,1.88) -- (2.62,1.88) -- cycle;

  % Hinge circle and mass m
  \draw[fill=white] (2.70,1.80) circle (0.11);
  \draw[fill=white] (5.40,1.80) circle (0.24);
  \node[above=2pt] at (5.40,2.04) {$m$};

  % Left guide slot and slider block
  \fill[pattern=north east lines] (-0.20,0.82) rectangle (0.65,0.98);
  \draw[line width=0.9pt] (-0.20,0.82) -- (0.65,0.82);
  \fill[pattern=north east lines] (-0.20,-0.18) rectangle (0.65,-0.02);
  \draw[line width=0.9pt] (-0.20,-0.02) -- (0.65,-0.02);
  \draw (0.04,0.08) rectangle (0.42,0.72);

  % Spring k between slider and vertical bar
  \draw (0.42,0.40) -- (1.02,0.40);
  \draw[decorate, decoration={coil, aspect=0.4, segment length=4.5pt, amplitude=3.5pt}] (1.02,0.40) -- (1.88,0.40);
  \draw (1.88,0.40) -- (2.62,0.40);
  \node[above=4pt] at (1.45,0.40) {$k$};

  % Dashpot gamma to the right wall
  \draw (2.78,0.40) -- (3.85,0.40);
  \draw (3.85,0.26) -- (3.85,0.54);
  \draw (3.65,0.18) -- (4.10,0.18) -- (4.10,0.62) -- (3.65,0.62);
  \draw (4.10,0.40) -- (5.05,0.40);
  \node[below=3pt] at (3.92,0.18) {$\gamma$};
  \fill[pattern=north east lines] (5.05,0.05) rectangle (5.22,0.75);
  \draw[line width=0.9pt] (5.05,0.05) -- (5.05,0.75);

  % Dimension l (vertical)
  \draw[thin] (2.15,1.80) -- (2.56,1.80);
  \draw[<->] (2.28,0.40) -- (2.28,1.80) node[midway, left=2pt] {$l$};

  % Dimension 2l (horizontal)
  \draw[thin] (5.40,1.52) -- (5.40,1.15);
  \draw[<->] (2.78,1.32) -- (5.40,1.32) node[midway, above=2pt] {$2l$};

  % Coordinate axis x_s at bottom left
  \draw[->] (-0.25,-0.55) -- (1.40,-0.55);
  \node[below=2pt] at (-0.25,-0.55) {$O$};
  \node[below=2pt] at (1.30,-0.55) {$x_s$};
\end{tikzpicture}

\end{center}
\end{problem}
\clearpage


\begin{problem}{3.2.15}
一个无阻尼的弹簧振子的受迫振动, 若强迫力的角频率 $\Omega = 5\mathrm{s}^{-1}$ 时发生共振, 给质量块增加 $1\mathrm{kg}$ 的质量, 当 $\Omega = 4.8\mathrm{s}^{-1}$ 时发生共振. 求弹簧振子原来的质量块的质量及弹簧的劲度系数.
\end{problem}
\clearpage


\begin{problem}{3.2.16}
如图 3.32 所示,一台自由运行的马达,为了减少其振动被放置在一块薄的橡皮垫上,马达埋入橡皮垫 10cm. 估算马达出现最大竖直振动时的转速(每分钟转数).
\begin{center}
\begin{tikzpicture}[>=Stealth, line width=0.7pt]
  % Rubber pad depression and hatching
  \fill[pattern=north east lines]
    (-2.05,0.42) .. controls (-1.60,0.62) and (-1.15,0.35) .. (-0.92,0)
    -- (0.92,0)
    .. controls (1.15,0.35) and (1.60,0.62) .. (2.05,0.42)
    -- (2.05,-0.25) -- (-2.05,-0.25) -- cycle;
  \draw[line width=0.8pt]
    (-2.05,0.42) .. controls (-1.60,0.62) and (-1.15,0.35) .. (-0.92,0)
    -- (0.92,0)
    .. controls (1.15,0.35) and (1.60,0.62) .. (2.05,0.42);

  % Motor pedestal and concentric circles
  \draw (-0.75,0) -- (-0.54,0.66);
  \draw (0.75,0) -- (0.54,0.66);
  \draw[fill=white] (0,1.42) circle (0.95);
  \draw (0,1.42) circle (0.48);

  % 10cm depth dimension on the right
  \draw[thin] (1.02,0) -- (2.55,0);
  \draw[thin] (1.72,0.50) -- (2.55,0.50);
  \draw[<->] (2.36,0) -- (2.36,0.50);
  \node[right=2pt, rotate=90, anchor=north] at (2.42,0.25) {$10\,\mathrm{cm}$};
\end{tikzpicture}

\end{center}
\end{problem}
\clearpage


\begin{problem}{3.2.17}
如图 3.33 所示, 一辆汽车沿 x 方向行驶并保持不变的水平速度 v, 汽车经过一个隆起部分, 其形状为

\[
y _ {0} = A \left(1 - \cos \frac {\pi x}{l}\right)
\]

$(0 \leqslant x \leqslant 2l)$ , 其他地方 $y_0 = 0$ . 求汽车在经过隆起部分时在竖直方向的运动. 将汽车简化为一个质量为 $m$ 的质点放在一个原长为 $l_0$ 、劲度系数为 $k$ 的弹簧上。忽略一切阻力并假定弹簧始终保持竖直。
\begin{center}
\begin{tikzpicture}[>=Stealth, line width=0.7pt]
  % Axes
  \draw[->] (-0.6,0) -- (5.1,0) node[below=2pt] {$x$};
  \draw[->] (0,-0.4) -- (0,2.2) node[left=2pt] {$y$};
  \node[below left=1pt] at (0,0) {$O$};

  % Bump curve: y0 = A(1 - cos(pi x / l)) on [0, 2l], with l = 2.1, A = 0.5
  \draw[line width=0.8pt] plot[domain=0:4.2, samples=80] (\x, {0.5*(1 - cos(deg(pi*\x/2.1)))});

  % Dashed vertical line at x = l and axis labels
  \draw[dashed] (2.1,0) -- (2.1,1.0);
  \node[below=2pt] at (2.1,0) {$l$};
  \node[below=2pt] at (4.2,0) {$2l$};

  % Left snapshot of mass-spring system
  \draw (0.30,0.02) -- (0.30,0.20);
  \draw[decorate, decoration={coil, aspect=0.4, segment length=3.2pt, amplitude=2.5pt}] (0.30,0.20) -- (0.30,0.68);
  \draw (0.30,0.68) -- (0.30,0.80);
  \draw[fill=white] (0.14,0.80) rectangle (0.46,1.02);
  \node[above=1pt] at (0.30,1.02) {$m$};
  \draw[->] (0.46,0.44) -- (0.95,0.44) node[right=1pt] {$v$};

  % Middle snapshot of mass-spring system
  \pgfmathsetmacro{\ymid}{0.5*(1 - cos(deg(pi*1.45/2.1)))}
  \draw (1.45,\ymid) -- (1.45,\ymid+0.18);
  \draw[decorate, decoration={coil, aspect=0.4, segment length=3.2pt, amplitude=2.5pt}] (1.45,\ymid+0.18) -- (1.45,\ymid+0.66);
  \draw (1.45,\ymid+0.66) -- (1.45,\ymid+0.78);
  \draw[fill=white] (1.29,\ymid+0.78) rectangle (1.61,\ymid+1.00);
  \draw[->] (1.61,\ymid+0.42) -- (2.10,\ymid+0.42) node[right=1pt] {$v$};
\end{tikzpicture}

\end{center}
\end{problem}
\clearpage


\begin{problem}{3.2.18}
一个由凸轮激励的、有阻尼的弹簧振子，如图3.34所示，凸轮使顶杆 $A$ 在铅垂方向按锯齿波规律做周期性运动，凸轮的升程为 $h$ ，周期为 $T$ ，顶杆又通过劲度系数为 $k_{1}$ 的弹簧使有阻尼的弹簧振子振动．弹簧振子的弹簧的劲度系数为 $k$ ，质量块质量为 $m$ ，阻尼系数为 $\gamma$ .求质量块的稳态振动.
\begin{center}
\begin{tikzpicture}[>=Stealth, line width=0.7pt]
  % Axes
  \draw[->] (0,0) -- (5.0,0) node[below=2pt] {$t$};
  \draw[->] (0,0) -- (0,2.1) node[left=2pt] {$y$};
  \node[below left=1pt] at (0,0) {$O$};

  % Tick for h
  \draw (0,1.4) -- (0.10,1.4);
  \node[left=2pt] at (0,1.4) {$h$};

  % Three sawtooth periods (T = 1.45)
  \draw[line width=0.8pt] (0,0) -- (1.45,1.4);
  \draw[dashed] (1.45,1.4) -- (1.45,0);
  \draw[line width=0.8pt] (1.45,0) -- (2.90,1.4);
  \draw[dashed] (2.90,1.4) -- (2.90,0);
  \draw[line width=0.8pt] (2.90,0) -- (4.35,1.4);
  \draw[dashed] (4.35,1.4) -- (4.35,0);

  % Period labels
  \node[below=2pt] at (1.45,0) {$T$};
  \node[below=2pt] at (2.90,0) {$2T$};
  \node[below=2pt] at (4.35,0) {$3T$};
\end{tikzpicture}

\end{center}
\begin{center}
\begin{tikzpicture}[>=Stealth, line width=0.7pt]
  % Top fixed support
  \fill[pattern=north east lines] (-1.25,4.85) rectangle (1.25,5.03);
  \draw[line width=0.9pt] (-1.25,4.85) -- (1.25,4.85);

  % Left branch: dashpot gamma
  \draw (-0.55,4.85) -- (-0.55,4.35);
  \draw (-0.77,3.90) -- (-0.77,4.35) -- (-0.33,4.35) -- (-0.33,3.90);
  \draw (-0.69,4.10) -- (-0.41,4.10);
  \draw (-0.55,4.10) -- (-0.55,3.30);
  \node[left=3pt] at (-0.77,4.12) {$\gamma$};

  % Right branch: spring k
  \draw (0.55,4.85) -- (0.55,4.52);
  \draw[decorate, decoration={coil, aspect=0.4, segment length=4.2pt, amplitude=3.5pt}] (0.55,4.52) -- (0.55,3.62);
  \draw (0.55,3.62) -- (0.55,3.30);
  \node[right=4pt] at (0.55,4.07) {$k$};

  % Mass block m
  \draw (-1.20,2.75) rectangle (1.20,3.30);
  \node at (0,3.025) {$m$};

  % Lower spring k1
  \draw (0,2.75) -- (0,2.38);
  \draw[decorate, decoration={coil, aspect=0.4, segment length=4.2pt, amplitude=3.5pt}] (0,2.38) -- (0,1.48);
  \draw (0,1.48) -- (0,1.18);
  \node[right=4pt] at (0,1.93) {$k_1$};

  % Guide walls and pushrod A
  \fill[pattern=north east lines] (-0.82,0.32) rectangle (-0.34,1.25);
  \draw (-0.82,1.25) -- (-0.34,1.25) -- (-0.34,0.32) -- (-0.82,0.32);
  \fill[pattern=north east lines] (0.34,0.32) rectangle (0.82,1.25);
  \draw (0.82,1.25) -- (0.34,1.25) -- (0.34,0.32) -- (0.82,0.32);
  \draw (-0.22,0.30) rectangle (0.22,1.18);
  \node at (0,0.74) {$A$};

  % Cam profile below pushrod A
  \draw[line width=0.8pt]
    (0,-0.80)
    .. controls (-0.65,-0.80) and (-0.72,-0.15) .. (-0.30,0.16)
    .. controls (-0.05,0.34) and (0.38,0.22) .. (0.62,0.02)
    .. controls (0.72,-0.07) and (0.52,-0.12) .. (0.45,-0.25)
    .. controls (0.40,-0.55) and (0.30,-0.80) .. (0,-0.80);

  % Cam pivot and fixed base
  \fill[pattern=north east lines] (-0.34,-0.62) rectangle (0.34,-0.50);
  \draw (-0.34,-0.50) -- (0.34,-0.50);
  \draw (-0.18,-0.50) -- (-0.06,-0.26);
  \draw (0.18,-0.50) -- (0.06,-0.26);
  \draw[fill=white] (0,-0.22) circle (0.08);

  % Rotation arrow on the left of cam
  \draw[->] (-0.78,-0.05) to[bend right=30] (-0.78,-0.75);
\end{tikzpicture}

\end{center}
\end{problem}
\clearpage


\begin{problem}{3.2.19}
(1)一个有阻尼的受迫谐振子,它的运动微分方程为

\[
m \ddot {x} = - m \omega_ {0} ^ {2} x - \gamma \dot {x} + A \cos \omega t
\]

它的平均能量损耗率多大？

(2) 一个非谐振子的运动微分方程为

\[
m \ddot {x} = - m \omega_ {0} ^ {2} x + \alpha x ^ {2} + A \cos \omega t
\]

其中 $\alpha$ 为一个小的常量, t=0 时, x=0, $\dot{x}=0$ , 求此后的运动到 $\alpha$ 的一次项.
\end{problem}
\clearpage


\begin{problem}{3.2.20}
无阻尼的非线性自由振动,质量为 m 的振子受力

\[
F (x) = - k x + m \lambda x ^ {2}
\]

其中 $\lambda$ 是小常量，求振子满足初始条件 t=0 时 $x(0)=A,\dot{x}(0)=0$ 的一级和二级近似解.
\end{problem}
\clearpage


\begin{problem}{3.2.21}
已知一个单摆的悬挂点做强迫竖直小振动 $\eta(t)=\eta_{0}\cos\omega t$ ，此单摆由长为 L 的无质量杆和杆端的质点组成。

(1) 假设单摆摆角振幅很小, $\eta_{0}\ll L$ ,求摆角 $\varphi$ 满足的运动微分方程;

(2) 求满足初始条件 t=0 时 $\varphi=a,\dot{\varphi}=0$ 的一级近似解；

(3) $w$ 为何值时发生共振, 写出共振时满足式(2)给的初始条件的一级近似解.
\end{problem}
\clearpage


\section{简谐波}

\begin{problem}{3.3.1}
一平面简谐波

\[
y = 0. 2 \sin [ \pi (0. 5 x - 3 0 t) ]
\]

x、y 的单位为厘米，t 的单位为秒．求波的振幅、波长、频率、波速以及 x=1cm 处质元振动的初相位．
\end{problem}
\clearpage


\begin{problem}{3.3.2}
一波的频率为 $20s^{-1}$ ，波速为 $80m \cdot s^{-1}$ ，振幅为 0.02m，求：

(1) 波的相位相差 $45^{\circ}$ 的两个点间的距离；

(2) 在一给定点处, 时间相隔 0.01s 的两位移之相位差;

(3) 在一给定点处, 时间相隔 0.01s 的质元的两个位置的最大距离;

(4) 在一给定点处, 质元的振动相位相差 $45^{\circ}$ 的质元的两个位置的最大距离.
\end{problem}
\clearpage


\begin{problem}{3.3.3}
一列沿长弦线向 $x$ 正方向传播的平面简谐波，位于 $x = x_{1} = 0$ 及 $x = x_{2} = 1\mathrm{m}$ 处两质元的简谐振动如下：

\[
y _ {1} = 0. 2 \sin 3 \pi t, y _ {2} = 0. 2 \sin \left(3 \pi t + \frac {\pi}{8}\right)
\]

$y$ 的单位为米， $t$ 的单位为秒．求：

(1) 频率；

(2) 波长 $\lambda;$ 

(3) 波速 v;

(4) 该两质元的相对速率的最大值.
\begin{center}
\begin{tikzpicture}[>=Stealth, line width=0.7pt]
  \def\L{4.5}
  \def\ang{22.5}
  \coordinate (O) at (0,0);
  \coordinate (A) at ({\L*cos(\ang)},{\L*sin(\ang)});
  \coordinate (B) at ({\L*cos(\ang)-\L},{\L*sin(\ang)});

  % Horizontal reference line
  \draw (O) -- (\L,0);

  % Vectors
  \draw[->, line width=0.8pt] (O) -- (A) node[midway, below right=2pt] {$0.6\pi$};
  \draw[->, line width=0.8pt] (A) -- (B) node[midway, above=2pt] {$0.6\pi$};
  \draw[->, line width=0.8pt] (O) -- (B) node[midway, left=3pt] {$|\dot{y}_2 - \dot{y}_1|_{\max}$};

  % Angle labels
  \node at (1.95,0.34) {$\dfrac{\pi}{8}$};
  \node at (2.45,1.35) {$\dfrac{\pi}{8}$};
\end{tikzpicture}

\end{center}
\end{problem}
\clearpage


\begin{problem}{3.3.4}
一平面简谐波以 $50\mathrm{m / s}$ 的速率沿 $x$ 负方向传播，在 $t = 1\mathrm{s}$ 时的波形如图3.37所示，写出此波的运动学方程，并画出 $t = 1\mathrm{s}$ 时各质元的 $\frac{\partial y}{\partial t} - x$ 图.
\begin{center}
\begin{tikzpicture}[>=Stealth, line width=0.7pt]
  % Axes
  \draw[->] (0,0) -- (5.3,0);
  \node[below=2pt] at (5.1,-0.05) {\shortstack{$x$\\[-2pt](m)}};
  \draw[->] (0,-1.55) -- (0,1.75);
  \node[left=2pt] at (-0.02,1.55) {\shortstack{$y$\\[-2pt](mm)}};

  % Wave curve: y(x) = 1.05 * sin(360 * x / 5.0 + 30)
  % Zero crossings at x = 2.0833 and x = 4.5833 (half-wavelength = 2.5 cm -> 10 m)
  \draw[line width=0.85pt] plot[domain=0:4.85, samples=100] (\x, {1.05*sin(360*\x/5.0 + 30)});

  % Dashed extrema lines and y-axis labels
  \draw[dashed] (0,1.05) -- (0.833,1.05);
  \draw[dashed] (0,-1.05) -- (3.333,-1.05);
  \node[left=2pt] at (0,1.05) {$10$};
  \node[left=2pt] at (0,0.525) {$5$};
  \node[left=2pt] at (0,0) {$0$};
  \node[left=2pt] at (0,-1.05) {$-10$};

  % Half-wavelength dimension 10
  \draw[thin] (2.0833,0.06) -- (2.0833,1.05);
  \draw[thin] (4.5833,0.06) -- (4.5833,1.05);
  \draw[<->] (2.0833,0.84) -- (4.5833,0.84) node[midway, above=2pt] {$10$};
\end{tikzpicture}

\end{center}
\end{problem}
\clearpage


\begin{problem}{3.3.5}
一根均匀的长弦线, 线密度为 0.02kg/m, 被 50N 的力拉紧, 弦线的一端 x=0

做简谐振动, 振幅为 0.02m, 周期为 1s, t=0 时, 位移 y 为 0.01m, $\frac{\partial y}{\partial t}$ 为负值. 写出向 x 正方向传播的简谐波的运动学方程.
\end{problem}
\clearpage


\begin{problem}{3.3.6}
一脉冲从长弦线的一端行进至另一端,需时0.1秒,弦线跨过一滑轮后悬一质量为弦线质量100倍的重物,问弦线的长度多大?
\end{problem}
\clearpage


\begin{problem}{3.3.7}
在拉紧的弦线上传播一个脉冲：

\[
y (x, t) = \frac {b ^ {3}}{b ^ {2} + (a x + u t) ^ {2}}
\]

其中 $a, b, u$ 均为正值常量，问：

(1) 此介质是色散介质吗？说明理由；

(2) 该脉冲速率多大？方向如何？

(3) 弦线上 x=0 处质元在 t=1 时的振动速度.
\end{problem}
\clearpage


\begin{problem}{3.3.8}
在两端固定、长度为 100m 的弦线上,有一如图 3.39 所示的脉冲以 40m/s 的速率向右行进,且保持其形状不变.

(1) 定性地画出在脉冲处于图示位置时弦线各质元的振动速度的轮廓图；

(2) 近似地估算弦线质元的最大振动速度；

(3) 若弦线总质量为 2kg, 弦线中的张力有多大?
\begin{center}
\begin{tikzpicture}[>=Stealth, line width=0.7pt]
  % Left and right fixed walls
  \fill[pattern=north east lines] (-0.16,-0.32) rectangle (0,0.32);
  \draw[line width=0.8pt] (0,-0.32) -- (0,0.32);
  \fill[pattern=north east lines] (6.00,-0.32) rectangle (6.16,0.32);
  \draw[line width=0.8pt] (6.00,-0.32) -- (6.00,0.32);

  % String baseline and triangular pulse
  \draw[line width=0.8pt] (0,0) -- (6.00,0);
  \draw[line width=0.8pt] (1.95,0) -- (2.40,0.36) -- (2.85,0);

  % Top dimensions (40m, 60m)
  \draw[thin] (0,0.52) -- (0,0.98);
  \draw[thin] (2.40,0.52) -- (2.40,0.98);
  \draw[thin] (6.00,0.52) -- (6.00,0.98);
  \draw[<->] (0,0.80) -- (2.40,0.80) node[midway, above=2pt] {$40\,\mathrm{m}$};
  \draw[<->] (2.40,0.80) -- (6.00,0.80) node[midway, above=2pt] {$60\,\mathrm{m}$};

  % Pulse height dimension (0.1m)
  \draw[thin] (2.82,0.36) -- (3.22,0.36);
  \draw[<->] (3.06,0) -- (3.06,0.36);
  \node[right=2pt] at (3.06,0.18) {$0.1\,\mathrm{m}$};

  % Pulse half-width dimensions (1m, 1m)
  \draw[thin] (1.95,-0.12) -- (1.95,-0.62);
  \draw[thin] (2.40,-0.12) -- (2.40,-0.62);
  \draw[thin] (2.85,-0.12) -- (2.85,-0.62);
  \draw[<->] (1.95,-0.46) -- (2.40,-0.46) node[midway, above=1pt] {$1\,\mathrm{m}$};
  \draw[<->] (2.40,-0.46) -- (2.85,-0.46) node[midway, above=1pt] {$1\,\mathrm{m}$};
\end{tikzpicture}

\end{center}
\end{problem}
\clearpage


\begin{problem}{3.3.9}
一个对称的三角形脉冲在弦线上沿 $x$ 正方向运动, 其最大高度为 $0.4 \mathrm{~m}$ , 全长为 $1.0 \mathrm{~m}$ , 波速为 $24 \mathrm{~m} / \mathrm{s}$ , 在 $t = 0$ 时, 整个脉冲位于 $x = 0$ 至 $x = 1 \mathrm{~m}$ 之间. 画出 $x = 1 \mathrm{~m}$ 处质元振动速度随时间变化的曲线.
\end{problem}
\clearpage


\begin{problem}{3.3.10}
气体或液体的压缩系数定义为 $k = -\frac{1}{\gamma} \frac{\mathrm{d}V}{\mathrm{d}p}$ , 这里 $-\mathrm{d}V$ 是压强增加 $\mathrm{d}p$ 时的体积减小量, 标准状态下空气的压缩系数比水约大 15000 倍.

(1) 推导声波速度公式 $v=\frac{1}{\sqrt{k\rho}}$ ，其中 $\rho$ 是质量

密度；

(2) 在标准状态下, 空气中的声速约为 $330 \mathrm{~m} / \mathrm{s}$ , 水中的声速约为 $1470 \mathrm{~m} / \mathrm{s}$ , 假设水是充满了微小气泡的混合物, 微小气泡仅占总体积的 $1 \%$ (忽略由气泡引起的混合物的密度和纯水的密度之间的差异), 求混合物的压缩系数 $k$ , 再求混合物中的声速 $v$ 值, 并与纯水或空气中的声速 $v$ 值相比较.
\end{problem}
\clearpage


\begin{problem}{3.3.11}
(1) 写出在弦中传播基波的频率与弦的物理性质、几何性质的关系；

(2) 从牛顿运动定律出发, 分析一小段弦所发生的变化导出上述关系.
\end{problem}
\clearpage


\begin{problem}{3.3.12}
声音在气体中的速度由下式计算: $v = \sqrt{\text{绝热体积模量/密度}}$ .

(1) 证明此公式的量纲是对的；

(2) 此公式隐含着认为声音在空气中的传播是准静态过程,但从另一方面,在空气分子的方均根速率大致为 500m/s 的温度下,空气中的声速大致为 340m/s,那么这种过程为什么能够是准静态的?
\end{problem}
\clearpage


\begin{problem}{3.3.13}
若有下述两列波在介质中传播：

\[
y _ {1} = A \sin (5 x - 1 0 t), y _ {2} = A \sin (4 x - 9 t)
\]

式中 x、y 以米为单位，t 以秒为单位。

(1) 写出合成波的方程；

(2) 求群速度 $v_{g}$ ;

(3) 求出合成波中振幅为零的两相邻点间的距离；

(4) 说明此介质是否是色散介质,说明理由.
\end{problem}
\clearpage


\begin{problem}{3.3.14}
两个不同的音叉在完全相同的两段长绳上产生稳定的简谐波, 振幅 $A_{1} = 2A_{2}$ , 波长 $\lambda_{1} = \frac{1}{2}\lambda_{2}$ . 设绳子除了与音叉交换能量外, 不与其他物体交换能量. 求两音叉给予绳子的平均功率之比.
\end{problem}
\clearpage


\begin{problem}{3.3.15}
一正弦式空气波沿直径为 0.14m 的圆柱形管子的轴向传播, 波的强度为 $9.0 \times 10^{-3} J/m^{2} \cdot s$ , 频率为 300Hz, 波速为 300m/s, 求:

(1) 波中平均能量密度和最大能量密度；

(2) 每两个相邻质元同振动相位面间的区域中含有的机械能.
\end{problem}
\clearpage


\begin{problem}{3.3.16}
一波源以 $3.5 \times 10^{4} \mathrm{~W}$ 的功率发射球面电磁波，在某处测得该波的平均能量密度为 $7.8 \times 10^{-15} \mathrm{~J} / \mathrm{m}^{3}$ . 求该处离波源的距离. 电磁波的传播速度为 $3.0 \times 10^{8} \mathrm{~m} / \mathrm{s}$ .
\end{problem}
\clearpage


\begin{problem}{3.3.17}
两列平面简谐波

\[
\begin{array}{r l} & y _ {1} = A _ {1} \cos \left[ \omega_ {1} \left(t - \frac {x}{v}\right) + \alpha_ {1} \right] \\ & y _ {2} = A _ {2} \cos \left[ \omega_ {2} \left(t - \frac {x}{v}\right) + \alpha_ {2} \right] \end{array}
\]

同时在同一介质中传播,分别讨论 $\omega_{1}\neq\omega_{2}$ 及 $\omega_{1}=\omega_{2}$ 两种情况下波强是否等于两波分别单独传播时的波强之和.
\end{problem}
\clearpage


\section{边界效应和干涉}

\begin{problem}{3.4.1}
设反射波的振动方向仍保持入射波的振动方向, 若介质的杨氏模量值为切变模量值的 3 倍, 试用惠更斯原理求出一列纵波的反射波是一列横波的入射角以及一列在入射面内振动的横波的反射波为一列纵波的入射角, (不考虑这样的入射波、反射波能否遵从波的叠加原理的问题).
\end{problem}
\clearpage


\begin{problem}{3.4.2}
有一稀奇的现象,当风从声源向观察者吹时,人们偶而听到来自远方的声音,惊人的清楚.

(1) 说明这个现象不能用“随风而来”来解释, 即均匀的风速不能说明这个现象;

(2) 吹遍地面的风, 有一垂直的速度梯度, 在地面附近可近似表示为 $v = ky^2$ , 这里 $y$ 是离地面的高度, $k$ 是常数. 对于给定的 $k$ 值和声速 $v_{s}$ , 计算顺风时出现声强最大增强处离声源的距离 $s$ ;

提示: 可以假定声波线遵循下列轨道: $y = h \sin \frac{\pi x}{s}$ .

（3）人们还注意到，甚至在无风时，声音在湖面上传播也有增强现象。在这种情况下，将发生什么？
\end{problem}
\clearpage


\begin{problem}{3.4.3}
当平面简谐波垂直入射于两种介质的界面时,证明:对于没有耗散和吸收的无色散介质,在任何时刻,入射波的波强等于反射波的波强和透射波的波强之和.
\end{problem}
\clearpage


\begin{problem}{3.4.4}
两根连在一起的弦线, 线密度分别为 $\eta_{1}$ 和 $\eta_{2}$ , 设想有一波从弦线 2 向弦线 1 传播, 设入射波方程为

\[
y _ {2} = A \cos (\omega t - k _ {2} x) \quad (x <   0)
\]

边界位于 $x = 0$ 处．试分别求出 $\frac{\eta_2}{\eta_1} = \frac{1}{4}$ 和 $\frac{\eta_2}{\eta_1} = 9$ 时的反射波、透射波的数学表达式.
\end{problem}
\clearpage


\begin{problem}{3.4.5}
一列波从线密度为 $\eta_{1}$ 的弦线传向线密度为 $\eta_{2}$ 的弦线，在连接点处反射波和透射波振幅之比 $r$ 、相位关系与 $\eta_{1}, \eta_{2}$ 间有何关系？
\end{problem}
\clearpage


\begin{problem}{3.4.6}
平面声波自空气中垂直入射到水面上, 求振幅反射系数 $\frac{B}{A}$ 和振幅透射系数 $\frac{C}{A}$ 

提示:空气和水的体积弹性模量不同,在界面处, $\frac{\partial(y_{1}+y'_{1})}{\partial x}=\frac{\partial y_{2}}{\partial x}$ 的边界条件需作修改.
\end{problem}
\clearpage


\begin{problem}{3.4.7}
一长列振幅为 $A$ 的波包, 主要由角频率很接近 $\omega_0$ 的波组成, 它在一根无限长的线密度为 $\eta$ 、张力为 $T$ 的弦上传播. 波包碰到弦上一个质量为 $m$ 的小珠, 设 $mg \ll T$ .

(1) 求透射波振幅；

(2) 在 $m, \omega_{0}$ 很大的极限下，透射波振幅与 $\omega_{0}$ 有何关系？
\begin{center}
\begin{tikzpicture}[>=Stealth, line width=0.7pt]
  % Horizontal string before and after wave packet
  \draw[line width=0.8pt] (0,0) -- (0.45,0);
  \draw[line width=0.8pt] plot[domain=0.45:2.61, samples=120] (\x, {0.26*sin(360*(\x-0.45)/0.72)});
  \draw[line width=0.8pt] (2.61,0) -- (5.50,0);

  % Bead of mass m
  \fill (4.10,0) circle (0.075);
  \node[above=3pt] at (4.10,0.06) {$m$};

  % Peak-to-peak amplitude 2A
  \draw[thin] (0.63,0.26) -- (1.03,0.26);
  \draw[thin] (0.63,-0.26) -- (1.03,-0.26);
  \draw[<->] (0.83,-0.26) -- (0.83,0.26);
  \draw[thin] (0.83,0) -- (1.38,0.52) node[above right=-2pt] {$2A$};
\end{tikzpicture}

\end{center}
\end{problem}
\clearpage


\begin{problem}{3.4.8}
一列简谐波沿线密度为 $\eta$ 、张力为 $T_0$ 的弦线从右入射，被套在竖直棒上的无质量小环反射，如图3.51所示，环受到的摩擦力的大小与速率成正比，比例系数为 $\gamma$ 。求振幅的反射系数 $\frac{B}{A}$ ，并讨论 $\gamma \to \infty$ 和 $\gamma \to 0$ 的两极限情况。
\begin{center}
\begin{tikzpicture}[>=Stealth]
  % Axes
  \draw[->] (-0.1, 0) -- (4.8, 0) node[below] {$x$};
  \draw[->] (0, -0.75) -- (0, 1.45) node[left] {$y$};
  \node[below left] at (0, 0) {$O$};

  % Wave profile and massless ring at x = 0
  \draw[thick, smooth] plot[domain=0:4.35, samples=80] (\x, {0.62*sin(deg(1.45*\x + 0.45))});
  \fill (0, {0.62*sin(deg(0.45))}) circle (1.8pt);
\end{tikzpicture}

\end{center}
\end{problem}
\clearpage


\begin{problem}{3.4.9}
若上题的小环具有质量 $m$ ，环与竖直棒间无摩擦，有一列角频率为 $\omega$ 的简谐波被小环反射，求振幅的反射系数 $\frac{B}{A}$ 及与入射波的相位差。（设 $mg \ll T_0$ ）.
\end{problem}
\clearpage


\begin{problem}{3.4.10}
线密度为 $\eta$ 的一根长弦处在张力 $T_{0}$ 下，一点质量 m 附着在弦上某一点，从

左边入射一角频率为 $\omega$ 的波沿这弦传播，设 $mg \ll T_{0}$ .

(1) 计算被点质量 m 反射的那部分能量与入射能量之比；

(2) 假定点质量 $m$ 被一线密度为 $\eta' (\gg \eta)$ 、长度 $l = \frac{m}{\eta'}$ 的短弦代替，固定 $m$ 值不变，变 $l$ 和 $\eta', l$ 满足什么条件，上述结果能近似成立？
\end{problem}
\clearpage


\begin{problem}{3.4.11}
一个点质量 $m$ 附在线密度为 $\eta$ 、张力为 $T$ 的无限长的弦上，同时和一个劲度系数为 $k$ 的弹簧相连，并受到一个外加的力 $F(t)$ ，如图3.52所示。点质量运动时，它在弦上产生一个起始波。

(1)求质点的运动微分方程和弦的运动方程；

(2) 若 $F(t)=F_{0}\cos\omega t$ ，求辐射到弦上的平均功率（不考虑瞬变过程）；

(3) 如果没有驱动力, 点质量的振动将衰减掉, 证明它与弦耦合引起的阻尼等同于一个摩擦阻力 $f = -\gamma v$ , 其中 $v$ 是点质量的速度, 求 $\gamma$ 与弦的参数之间的关系.
\begin{center}
\begin{tikzpicture}[>=Stealth]
  % Fixed ceiling with hatching
  \draw[thick] (-0.65, 2.15) -- (0.65, 2.15);
  \foreach \x in {-0.55, -0.35, -0.15, 0.05, 0.25, 0.45} {
    \draw (\x, 2.15) -- (\x+0.15, 2.32);
  }

  % Axes and external force F(t)
  \draw[->] (-2.7, 0) -- (2.8, 0) node[below] {$x$};
  \draw[->] (0, 2.15) -- (0, 2.65) node[left] {$y$};
  \draw[->, thick] (0, -0.85) -- (0, 0.22);
  \node[right=2pt] at (0, -0.45) {$F(t)$};
  \node[below left] at (0, 0) {$O$};

  % Spring and point mass m
  \draw[decorate, decoration={zigzag, segment length=5.5pt, amplitude=3pt, pre length=12pt, post length=12pt}]
    (0, 2.15) -- (0, 0.30);
  \node[right=4pt] at (0, 1.25) {$k$};
  \fill (0, 0.30) circle (2pt) node[above right=1pt] {$m$};

  % Symmetric wave along the string
  \draw[thick] (0, 0.30)
    .. controls (0.35, 0.16) and (0.70, 0.16) .. (1.05, 0.40)
    .. controls (1.35, 0.62) and (1.65, 0.66) .. (2.00, 0.56)
    .. controls (2.30, 0.46) and (2.48, 0.32) .. (2.60, 0.20);
  \draw[thick] (0, 0.30)
    .. controls (-0.35, 0.16) and (-0.70, 0.16) .. (-1.05, 0.40)
    .. controls (-1.35, 0.62) and (-1.65, 0.66) .. (-2.00, 0.56)
    .. controls (-2.30, 0.46) and (-2.48, 0.32) .. (-2.60, 0.20);

  % Transverse displacement markers y_1 and y_2
  \draw[dashed] (-1.60, 0) -- (-1.60, 0.63);
  \node[right=1pt] at (-1.60, 0.30) {$y_1$};
  \draw[dashed] (1.60, 0) -- (1.60, 0.63);
  \node[right=1pt] at (1.60, 0.30) {$y_2$};
\end{tikzpicture}

\end{center}
\end{problem}
\clearpage


\begin{problem}{3.4.12}
由劲度系数为 $k$ 、原长为 $a$ 的非常多相同的弹簧连结的非常长的珠子链，如图3.53所示，每个珠子沿 $x$ 方向做无阻尼振动，除一个珠子质量为 $m_0$ 外，其余的质量均为 $m (> m_0)$ ，不计弹簧质量。

(1)求在远离那个特殊的珠子处传播的波的波矢和角频率之间的关系；

(2) 当一个波矢为 k 的波碰到这个特殊珠子时, 反射概率多大?
\begin{center}
\begin{tikzpicture}[>=Stealth]
  % Broken spring ends for infinite chain
  \draw[decorate, decoration={zigzag, segment length=3pt, amplitude=1.8pt}] (-2.85, 0) -- (-2.48, 0);
  \draw[decorate, decoration={zigzag, segment length=3pt, amplitude=1.8pt}] (2.48, 0) -- (2.85, 0);

  % Springs between masses
  \foreach \i in {-4, -3, -2, 1, 2, 3} {
    \draw[decorate, decoration={zigzag, segment length=3pt, amplitude=1.8pt, pre length=2pt, post length=2pt}]
      (\i*0.62 + 0.055, 0) -- (\i*0.62 + 0.62 - 0.055, 0);
    \node[above=5pt] at (\i*0.62 + 0.31, 0) {\small $K$};
  }
  \draw[decorate, decoration={zigzag, segment length=3pt, amplitude=1.8pt, pre length=2pt, post length=2pt}]
    (-0.62 + 0.055, 0) -- (-0.13, 0);
  \node[above=5pt] at (-0.31, 0) {\small $K$};
  \draw[decorate, decoration={zigzag, segment length=3pt, amplitude=1.8pt, pre length=2pt, post length=2pt}]
    (0.13, 0) -- (0.62 - 0.055, 0);
  \node[above=5pt] at (0.31, 0) {\small $K$};

  % Identical beads m
  \foreach \i in {-4, -3, -2, -1, 1, 2, 3, 4} {
    \fill (\i*0.62, 0) circle (1.8pt) node[below=3pt] {$m$};
    \node[above=14pt] at (\i*0.62, 0) {\small $\i$};
  }

  % Special mass m_0 at index 0
  \draw[fill=white, thick] (-0.13, -0.13) rectangle (0.13, 0.13);
  \node[below=5pt] at (0, -0.10) {$m_0$};
  \node[above=14pt] at (0, 0) {\small $0$};

  % x-direction arrow
  \draw[->] (1.65, -0.48) -- (2.90, -0.48) node[below] {$x$};
\end{tikzpicture}

\end{center}
\end{problem}
\clearpage


\begin{problem}{3.4.13}
一质量为 $m$ 、原长为 $l$ 的弹簧, 其劲度系数为 $k$ . 弹簧的一端固定, 另一端系一质量为 $M$ 的物体, 水平放置. $M$ 在光滑的水平面上运动.

(1) 导出该系统纵向振动的波动方程；

(2) 求出作为 $m$ 的函数的最小简正频率。这里 $M$ 和 $k$ 都是有限的，且 $m \ll M$ 。
\end{problem}
\clearpage


\begin{problem}{3.4.14}
一根无限长的细绳, 线密度为 $\eta$ , 张力为 $T$ . 在 $t = 0$ 时, 绳偏离平衡位置的位移为 $f(x)$ . 初始速度为 $g(x)$ , 求 $t > 0$ 时绳子的运动.
\end{problem}
\clearpage


\begin{problem}{3.4.15}
一根线密度为 $\eta$ 、张力为 $T$ 的绳， $x = 0$ 端固定， $x = l$ 端终止在一个不计质量的环上，环能在一垂直棒上无摩擦自由滑动、写 $x = l$ 端的边界条件，并求振动的简正模。
\end{problem}
\clearpage


\begin{problem}{3.4.16}
若上题环的质量为 $m$ ，求振动的简正模，并证明： $m \to \infty$ 时，简正模趋于两端固定的情况。
\end{problem}
\clearpage


\begin{problem}{3.4.17}
两种线密度的弦线在 $x = 0$ 处连接， $x > 0$ 处， $\eta = \eta_{2}, x < 0$ 处， $\eta = \eta_{1}$ ，张力为 $T$ 。一列简谐波从弦线1入射到连接点。分别 $\eta_{1} > \eta_{2}$ 和 $\eta_{1} < \eta_{2}$ 两种情况，讨论反射波、透

\[
\begin{array}{r l} M \frac {\mathrm{d} ^ {2} y _ {n}}{\mathrm{d} t ^ {2}} & = T \left(\frac {y _ {n + 1} - y _ {n}}{a}\right) - T \left(\frac {y _ {n} - y _ {n - 1}}{a}\right) \\ & = \frac {T}{a} (y _ {n + 1} - 2 y _ {n} + y _ {n - 1}) \quad n = 1, 2, \dots , N \end{array}
\]

其中 $y_0 = y_{n + 1} = 0$ 

由于 $N$ 相当大, 要用第十章求少自由度的小振动的简正模式的办法, 解 $N \times N$ 的行列式是困难的. 可用求解3.4.12的办法.

在同一振动模式中,各珠子有同样的角频率,相位

相同或相反,振幅不同,其解可写成

\[
y _ {n} = A _ {n} \cos (\omega t + \varphi) \quad n = 1, 2, \dots , N
\]

参照两端固定的弦的振动模式是各种驻波,对此珠子链,试用

\[
A _ {n} = A \sin k n a + B \cos k n a \quad n = 0, 1, \dots , N + 1
\]

由边界条件 $A_0 = 0$ 得 $B = 0$ ，故

\[
A _ {n} = A \sin k n a
\]

$A_{N + 1} = 0, \sin k(N + 1)a = 0$ ，即 $\sin kL = 0$ 

\[
k _ {m} l = m \pi ,
\]

\[
k _ {m} = \frac {m \pi}{L}, m = 1, 2, \dots , N
\]

将

$y_{n} = A\sin knacos(\omega t + \varphi)$ 代入方程组，得

\[
M \omega^ {2} \sin k n a = \frac {T}{a} [ \sin k (n + 1) a - 2 \sin k n a + \sin k (n - 1) a ]
\]

用三角公式 $\sin \alpha +\sin \beta = 2\sin \frac{\alpha + \beta}{2}\cos \frac{\alpha - \beta}{2}$ 

上式化为

\[
- M \omega^ {2} = \frac {2 T}{a} (\cos k a - 1)
\]

故

\[
\omega_ {m} ^ {2} = \frac {2 T}{M a} \left(1 - \cos \frac {m \pi a}{L}\right) m = 1, 2, \dots , N
\]

作几点附带说明：

\[
1) y _ {n} = A \sin k _ {m} n a \cos (\omega_ {m} t + \varphi_ {m})
\]

用 $\sin \alpha \cos \beta = \frac{1}{2} [\sin (\alpha +\beta) + \sin (\alpha -\beta)]$ 可改写为

\[
y _ {n} = \frac {1}{2} A \left[ \sin \left(k _ {m} n a + \omega_ {m} t + \varphi_ {m}\right) + \sin \left(k _ {m} n a - \omega_ {m} t - \varphi_ {m}\right) \right]
\]

第一项表示向左边传播的波,第二项表示向右边传播的波.

波的传播速度有两种计算方法:考虑向右传播的波

\[
\sin (k _ {m} n a - \omega_ {m} t - \varphi_ {n})
\]

考虑相邻的两个珠子的振动在紧挨的处于同相位的两个时刻

\[
\begin{array}{r l} k _ {m} n a - \omega_ {m} t _ {1} - \varphi_ {m} & = k _ {m} (n + 1) a - \omega_ {m} t _ {2} - \varphi_ {m} \\ v _ {m} & = \frac {a}{t _ {2} - t _ {1}} = \frac {\omega_ {m}}{k _ {m}} \end{array}
\]

令 $\omega_{m}^{2} = \frac{2T}{Ma} (1 - \cos k_{m}a),\frac{\omega_{m}}{k_{m}}\neq$ 常量(对于不同的 $\omega_{m}$ ).说明对于不同角频率的振动，其相

位的传播速度是不同的.

波速的另一种计算方法是把 $y_{n}$ 改成 $y(x)$ 的形式： $\sin (k_m n a - \omega_m t - \varphi_m)$ 改成 $\sin (k_m x - \omega_m t - \varphi_m) = -\sin \left[\omega_m \left(t - \frac{k_m}{\omega_m} x\right) + \varphi_m\right]$ ，即可看出 $v_m = \frac{\omega_m}{k_m}$ .

2) 考虑连续的极限情况．保持 L 不变，令 $N \to \infty, a \to 0, \frac{M}{a} = \eta$ 

\[
\begin{array}{r l}\omega_ {m} ^ {2}&= \frac {2 T}{M a} (1 - \cos k _ {m} a) = \frac {4 T}{M a} \sin^ {2} \frac {1}{2} k _ {m} a\\&= \frac {4 T}{M a} \cdot \left(\frac {1}{2} k _ {m} a\right) ^ {2} \left(\frac {\sin \frac {1}{2} k _ {m} a}{\frac {1}{2} k _ {m} a}\right) ^ {2}\\&\rightarrow \frac {4 T}{M a} \cdot \frac {1}{4} k _ {m} a ^ {2} = \frac {T}{\eta} k _ {m} ^ {2} = \frac {T}{\eta} \frac {m ^ {2} \pi^ {2}}{L ^ {2}}\\&\omega_ {m} \rightarrow \frac {m \pi}{L} \sqrt {\frac {T}{\eta}}\end{array}
\]

与两端固定的连续弦的情况完全相同.

3)对于有 N 个珠子的珠子链, $k_{m}=\frac{m\pi}{L}$ , m 可取 1,2, $\cdots$ , N,为什么不再往下取值? 这是因为这个系统具有 N 个自由度,由第十章讲的小振动理论可知,它至多有 N 个简正频率.这里也可以看到,当 $n=N+1$ 时,L 等于 $N+1$ 个半波长,所有珠子在此“振动模式”中都保持不动,这是没有意义的.
\begin{center}
\begin{tikzpicture}[>=Stealth]
  % Axes
  \draw[->] (-0.1, 0) -- (4.5, 0) node[below] {$x$};
  \draw[->] (0, -0.1) -- (0, 2.55) node[left] {$y$};

  % Coordinate projections
  \draw[dashed] (1.2, 0) node[below] {$(n-1)a$} -- (1.2, 0.72) -- (0, 0.72) node[left] {$y_{n-1}$};
  \draw[dashed] (2.35, 0) node[below] {$na$} -- (2.35, 1.25) -- (0, 1.25) node[left] {$y_n$};
  \draw[dashed] (3.5, 0) node[below] {$(n+1)a$} -- (3.5, 1.82) -- (0, 1.82) node[left] {$y_{n+1}$};

  % String segments with tension direction arrows
  \draw[thick] (0.55, 0.92) -- (1.2, 0.72);
  \draw[thick, postaction={decorate}, decoration={markings, mark=at position 0.52 with {\arrow{Stealth}}}]
    (2.35, 1.25) -- (1.2, 0.72);
  \draw[thick, postaction={decorate}, decoration={markings, mark=at position 0.52 with {\arrow{Stealth}}}]
    (2.35, 1.25) -- (3.5, 1.82);
  \draw[thick] (3.5, 1.82) -- (4.15, 2.18);

  % Mass points and labels
  \fill (1.2, 0.72) circle (1.8pt);
  \node[above=3pt] at (1.2, 0.76) {$T_{n-1}$};
  \fill (2.35, 1.25) circle (1.8pt);
  \node[above=3pt] at (2.35, 1.25) {$M$};
  \fill (3.5, 1.82) circle (1.8pt);
  \node[above=3pt] at (3.5, 1.82) {$T_{n+1}$};
\end{tikzpicture}

\end{center}
\end{problem}
\clearpage


\begin{problem}{3.4.19}
$S_{1}, S_{2}$ 是两相干波源，相距 $\frac{1}{4}$ 波长， $S_{1}$ 比 $S_{2}$ 相位超前 $\frac{\pi}{2}$ . 设两波在 $S_{1}, S_{2}$ 连线上的强度相同，均为 $I_{0}$ ，且不随距离变化，问 $S_{1}, S_{2}$ 连线上在 $S_{1}$ 外侧各点处的强度 $I_{1}$ 多大？在 $S_{2}$ 外侧各点处的强度 $I_{2}$ 多大？
\begin{center}
\begin{tikzpicture}[>=Stealth]
  % Source line and source ticks
  \draw[thick] (-0.25, 1.2) -- (4.2, 1.2);
  \draw[thick] (0, 1.2) -- (0, 1.35) node[above] {$S_1$};
  \draw[thick] (2.5, 1.2) -- (2.5, 1.35) node[above] {$S_2$};

  % Distance dimension (1/4 lambda)
  \draw (0, 1.05) -- (0, 0.50);
  \draw (2.5, 1.05) -- (2.5, 0.50);
  \draw[<->] (0, 0.78) -- (2.5, 0.78) node[midway, fill=white, inner sep=2pt] {$\frac{1}{4}\lambda$};

  % x-axis
  \draw[->, thick] (0, 0) -- (4.6, 0) node[right] {$x$};
  \node[below] at (0, 0) {$O$};
\end{tikzpicture}

\end{center}
\end{problem}
\clearpage


\begin{problem}{3.4.20}
$S_{1}, S_{2}$ 是两相干波源，相距 $30\mathrm{m}$ ， $S_{1}, S_{2}$ 的相位相反，频率 $\nu = 100\mathrm{Hz}$ ，波速 $\nu = 400\mathrm{m/s}$ ，在 $S_{1}, S_{2}$ 连线上两波振幅相同，不随距离变化，求 $S_{1}, S_{2}$ 连线上因干涉而静止的点的位置。
\end{problem}
\clearpage


\begin{problem}{3.4.21}
图中 $O$ 处有一波源，沿 $x$ 轴向两边发射振幅为 $A$ 、角频率为 $\omega$ 的简谐波，波速为 $v$ . $BC$ 为波密介质的反射面，入射波在此被完全反射， $BC$ 位于 $x = -d$ 处，若 $d = \frac{5}{4}\lambda, \lambda$ 为波长，求出 $x > -d$ 各处波动的数学表达式。设波源的初相位为零。
\begin{center}
\begin{tikzpicture}[>=Stealth]
  % Reflecting wall BC with hatching on the left
  \draw[thick] (0, -0.85) -- (0, 1.55);
  \foreach \y in {-0.8, -0.6, -0.4, -0.2, 0.0, 0.2, 0.4, 0.6, 0.8, 1.0, 1.2, 1.4} {
    \draw (0, \y) -- (-0.22, \y - 0.22);
  }
  \node[right=2pt] at (0, 1.45) {$B$};
  \node[right=2pt] at (0, -0.75) {$C$};

  % x-axis and origin O
  \draw[->, thick] (0, 0) -- (4.3, 0) node[below] {$x$};
  \draw (2.1, 0) -- (2.1, 0.08);
  \node[below=2pt] at (2.1, 0) {$O$};

  % Distance d between wall and O
  \draw (2.1, 0.45) -- (2.1, 0.95);
  \draw[<->] (0, 0.70) -- (2.1, 0.70) node[midway, above=2pt] {$d$};
\end{tikzpicture}

\end{center}
\end{problem}
\clearpage


\begin{problem}{3.4.22}
在弦线上传播的波,其表达式为

\[
y = 3 \cos \left[ 2 \pi \left(\frac {t}{0 . 1} - \frac {x}{1 0}\right) - \frac {\pi}{2} \right]
\]

为在弦线上形成驻波，在 $x = 1$ 处为波节

(1) 写出应叠加的波的表达式；

(2) 写出形成的驻波的表达式；

(3) 若弦线的线密度为 $1.0\mathrm{g/cm}$ , 求相邻两波节之间的总能量。所给的表达式中, $x$ 、 $y$ 均以米为单位, $t$ 以秒为单位。
\end{problem}
\clearpage


\begin{problem}{3.4.23}
在位于 $x$ 轴的弦线上有一驻波，测得 $x = n + \frac{5}{6} (\mathrm{m}) (n = 0, \pm 1, \pm 2, \dots)$ 处为波节，在波腹处，最大位移 $y_{\max}$ 为 $5\mathrm{m}$ ，从平衡位置到最大位移历时 $0.5\mathrm{s}$ ，以弦线上所有质元均处于平衡位置时开始计时，写出该驻波的表达式。
\end{problem}
\clearpage


\begin{problem}{3.4.24}
设平面横波 1 沿 BP 方向传播, 它在 B 点的振动方程为

\[
y _ {1} = 0. 2 \times 1 0 ^ {- 2} \cos 2 \pi t (\mathrm{m})
\]

平面横波 2 沿 CP 方向传播, 它在 C 点的振动方程为

\[
y _ {2} = 0. 2 \times 1 0 ^ {- 2} \cos (2 \pi t + \pi) (\mathrm{m})
\]

t 以秒计 .P 与 B、C 分别相距 0.40m 和 0.50m, BP、CP 的夹角 $\theta=60^{\circ}$ ，波速为 $0.20ms^{-1}$ ，求：

(1) 两波传到 P 处时引起质元振动的相位差；

(2) 分别考虑横波的振动方向均垂直于纸面和均平行于纸面两种情况, P 处质元合振动的振幅.
\begin{center}
\begin{tikzpicture}[>=Stealth]
  % Propagation rays intersecting at P
  \draw[->, thick] (-2.1, 1.21) -- (2.1, -1.21);
  \draw[thick] (-2.1, -1.21) -- (1.65, 0.95);

  % Labels
  \node[left=2pt] at (-2.1, 1.21) {$B$};
  \node[left=2pt] at (-2.1, -1.21) {$C$};
  \node[below=3pt] at (0, 0) {$P$};
  \node at (-0.65, 0) {$\theta$};
\end{tikzpicture}

\end{center}
\end{problem}
\clearpage


\begin{problem}{3.4.25}
图 3.61 中画的是演示声波的干涉的仪器的示意图, S 为声源, D 为耳朵或其他声音探测器, 路径 SAD 的长度是固定的, 路径 SBD 的长度可以变化. 干涉仪内是空气, 空气中声速为 $330 \, ms^{-1}$ . 现测得声音强度在 B 的第一位置时为极小值 100 单位, 渐增至 B 距第一位置为 $1.65 \times 10^{-2} \, m$ 的第二位置时, 有极大值 900 单位, 求:

(1)声波频率；

(2) 抵达探测器 D 经 A 和经 B 的两波的振幅之比(SBD 路程长于 SAD);

(3) 解释为什么同一声源发出的这两个波会有不同的振幅.
\begin{center}
\begin{tikzpicture}[>=Stealth]
  % Upper fixed U-tube (path A) with left/right ports
  \draw[thick] (-1.28, -1.15) -- (-1.28, -0.20) -- (-1.78, -0.20);
  \draw[thick] (-1.78, 0.20) -- (-1.28, 0.20) -- (-1.28, 1.0) arc (180:0:1.28)
    -- (1.28, 0.20) -- (1.78, 0.20);
  \draw[thick] (1.78, -0.20) -- (1.28, -0.20) -- (1.28, -1.15);
  \draw[thick] (-0.78, -1.15) -- (-0.78, 1.0) arc (180:0:0.78) -- (0.78, -1.15);
  \node[below=4pt] at (0, 1.78) {$A$};

  % Lower sliding U-tube (path B)
  \draw[thick] (-0.70, -0.82) -- (-0.70, -1.30) arc (180:360:0.70) -- (0.70, -0.82);
  \draw[thick] (-1.36, -0.82) -- (-1.36, -1.30) arc (180:360:1.36) -- (1.36, -0.82);
  \node[above=4pt] at (0, -2.00) {$B$};

  % Sound source S (speaker icon)
  \draw[thick] (-2.28, -0.10) rectangle (-2.12, 0.10);
  \draw[thick] (-2.12, 0.10) -- (-1.92, 0.22) -- (-1.92, -0.22) -- (-2.12, -0.10);
  \node[below=4pt] at (-2.12, -0.22) {$S$};

  % Detector D (ear icon)
  \draw[thick] (2.00, 0.12)
    .. controls (2.05, 0.24) and (2.22, 0.22) .. (2.22, 0.04)
    .. controls (2.22, -0.12) and (2.10, -0.24) .. (2.02, -0.16)
    .. controls (1.98, -0.11) and (2.02, -0.05) .. (2.06, -0.06);
  \draw[thick] (2.06, 0.06) .. controls (2.13, 0.10) and (2.15, -0.02) .. (2.08, -0.04);
  \node[below=4pt] at (2.10, -0.22) {$D$};
\end{tikzpicture}

\end{center}
\end{problem}
\clearpage


\begin{problem}{3.4.26}
图 3.62 中所示的是某一时刻入射波的波形图, 分别就 BC 是波密介质和波疏介质的反射面的两种情况, 画出此刻反射波波形的示意图(注意未说做完全反射, 反射波振幅一般应小于入射波振幅).
\begin{center}
\begin{tikzpicture}[>=Stealth]
  % Reflecting wall BC with hatching on the right
  \draw[thick] (3.8, -1.6) -- (3.8, 1.6);
  \foreach \y in {-1.5, -1.3, -1.1, -0.9, -0.7, -0.5, -0.3, -0.1, 0.1, 0.3, 0.5, 0.7, 0.9, 1.1, 1.3, 1.5} {
    \draw (3.8, \y) -- (4.05, \y + 0.25);
  }
  \node[left=2pt] at (3.8, 1.45) {$B$};
  \node[left=2pt] at (3.8, -1.45) {$C$};

  % Equilibrium axis
  \draw[thick] (-0.3, 0) -- (4.3, 0);

  % Incident wave profile
  \draw[thick, smooth] plot[domain=0.05:3.8, samples=80] (\x, {1.25*sin(deg(1.48*(\x - 0.20)))});
\end{tikzpicture}

\end{center}
\end{problem}
\clearpage


\begin{problem}{3.4.27}
一根线密度为 $0.15\mathrm{g / cm}$ 的弦线，其一端与频率为 $50\mathrm{Hz}$ 的音叉相连，另一端跨过一定滑轮后悬一质量为 $M$ 的重物，音叉到滑轮的距离为 $1\mathrm{m}$ 。音叉振动时，为使弦线上形成一个、二个、三个波腹，重物的质量 $M$ 应各取何值？音叉振动的振幅很小，近似认为处于波节。
\end{problem}
\clearpage


\begin{problem}{3.4.28}
两个同频率、同响度的音叉在某一点 A 处是节点, 即两音叉同时发声, 在 A 点听到的强度为零, 可任何一个音叉单独发声时, 在 A 点听到的强度都为 $I_{0}$ , 这与能量守恒定律有矛盾吗?
\end{problem}
\clearpage


\begin{problem}{3.4.29}
长 $l$ 、线密度 $\eta$ 的均质弦线两端固定，在强大张力 $T_0$ 作用下处于直线平衡位置，在离一端距离 $a$ 处作用一个简谐力 $F = F_0\cos \omega t$ ，其中 $F_0, \omega$ 为常量，若 $\sin \left(\omega \sqrt{\frac{\eta}{T_0} l}\right) \neq 0$ ，求弦线的运动学方程。
\end{problem}
\clearpage


\begin{problem}{3.4.30}
试用分离变量法作为受迫振动重解上题,并证明所得结果是一致的.
\end{problem}
\clearpage


\begin{problem}{3.4.31}
一作用于单位长度的强迫力 $f(x,t) = f_0\sin \left(\frac{n\pi x}{l}\right)\cos \omega t$ (其中 $f_0, n, \omega$ 均为常量, $n$ 为正整数)加于两端固定的、长 $l$ 、张力为一、线密度为 $\eta$ 的绳上, 绳子受到与速度成正比的阻力, 单位长度绳所受阻力的阻尼系数为 $\gamma$ , 求绳子的稳态运动.
\end{problem}
\clearpage


\begin{problem}{3.4.32}
一均匀细绳长度为 $l$ ，单位长度质量为 $\eta$ ，在平面 $(x, y)$ 内做横向小振动，它的端点 $(0, 0)$ 和 $(l, 0)$ 分别固定，张力为 $\tau$ ，振动时，受到与速度有关的摩擦力，如微元长 $\delta l$ ，有横向速度 $v$ ，则摩擦力为 $-\gamma v \delta l$ ，小振动 $y(x, t)$ 满足下列偏微分方程和边界条件：

\[
\begin{array}{c} \frac {\partial^ {2} y}{\partial t ^ {2}} + a \frac {\partial y}{\partial t} = b \frac {\partial^ {2} y}{\partial x ^ {2}} \\ y (0, t) = 0, y (l, t) = 0 \end{array}
\]

(1) 写出常量 a 和 b;

(2) 假定 $a^2 < \frac{b}{l^2}$ , 求 $y(x, t)$ (通解);

(3) 设 $y(x,0) = 0$ , $\dot{y}(x,0) = A\sin (3\pi x / l) + B\sin (5\pi x / l)$ , 其中 $A, B$ 为常量, 求 $y(x,t)$ ;

(4) 假设改为 $a=0,\dot{y}(x,0)=0$ ,

\[
y (x, 0) = \left\{ \begin{array}{c c} A x & 0 \leqslant x \leqslant \frac {1}{2} l \\ A (l - x) & \frac {1}{2} l \leqslant x \leqslant l \end{array} \right.
\]

求 $y(x,t)$
\end{problem}
\clearpage


\begin{problem}{3.4.33}
一条长 $l$ 、线密度为 $\eta$ 的软绳，上端固定，下端自由，在一个平面内做微小横振动。

(1) 导出偏微分方程, 并由此导出关于振动模式的微分方程;

(2) 用幂级数解法求解此微分方程(不把方程变为贝塞尔方程求解), 然后用数值解法近似算出最低简正频率.
\end{problem}
\clearpage


\begin{problem}{3.4.34}
一根长度为 l、线密度为 $\eta$ 的细绳以角速度 $\Omega$ 绕一根垂直的轴在水平面内转动，同时细绳在竖直方向做小振动，如图 3.67 所示．求简正频率（重力可忽略不计）.

提示: 勒让德多项式 $p_{l}(\cos\theta)$ 满足的方程为

\[
\frac {\mathrm{d}}{\mathrm{d} \cos \theta} \left(\sin^ {2} \theta \frac {\mathrm{d} p _ {l}}{\mathrm{d} \cos \theta}\right) = - l (l + 1) p _ {l} \quad l = 0, 1, 2, \dots
\]
\begin{center}
\begin{tikzpicture}[>=Stealth]
  % Axes
  \draw[->] (-0.1, 0) -- (4.5, 0) node[below] {$x$};
  \draw[->] (0, -0.85) -- (0, 1.85) node[left] {$y$};

  % Angular velocity Omega around y-axis
  \draw[->, thick] (-0.21, 1.41) arc (145:395:0.26 and 0.10);
  \node[right=3pt] at (0.24, 1.42) {$\Omega$};

  % Vibrating rotating string
  \draw[thick] plot[smooth] coordinates {
    (0, 0) (0.55, 0.22) (1.10, 0) (1.70, -0.16) (2.30, 0)
    (2.85, 0.22) (3.35, 0) (3.68, -0.17) (3.92, 0) (4.05, 0.18)
  };

  % String length L
  \draw (4.05, -0.28) -- (4.05, -0.75);
  \draw[<->] (0, -0.55) -- (4.05, -0.55) node[midway, above=2pt] {$L$};
\end{tikzpicture}

\end{center}
\end{problem}
\clearpage


\section{声波}

\begin{problem}{3.5.1}
一开管风琴管(两端开口)的基频为 300Hz, 一闭管风琴管(一端封闭)的第二泛音和该开管风琴管的第二泛音频率相同, 问两根管子的长度各是多少? 空气中声速为 330m/s.
\end{problem}
\clearpage


\begin{problem}{3.5.2}
一小提琴的弦的质量为 $2.0\mathrm{g}$ , 它的两固定点之间的距离为 $50\mathrm{cm}$ , 不按手指演奏时发出的声音是 $A$ 调 $440\mathrm{Hz}$ . 试问要奏出 $C$ 调 $528\mathrm{Hz}$ , 手指应按在什么位置? 弦线张力多大?
\end{problem}
\clearpage


\begin{problem}{3.5.3}
有一口井,侧面是竖直的,井中有水,它可与 $7\mathrm{Hz}$ 和 $7\mathrm{Hz}$ 以上的某些频率发生共鸣,井内的空气密度为 $1.1\mathrm{kg} / \mathrm{m}^3$ ,压强为 $9.5\times 10^{4}\mathrm{Pa}$ ,比热容比为 $\frac{7}{5}$ .试问这井水的水面有多深?比 $7\mathrm{Hz}$ 高的多大频率也能发生共鸣?
\end{problem}
\clearpage


\begin{problem}{3.5.4}
一小提琴的弦长 $31.6\mathrm{cm}$ , 线密度为 $0.65\mathrm{g/m}$ , 将此弦放在扬声器附近, 扬声器由一频率可变的音频振荡器来策动. 当振荡器的频率从 $500\mathrm{Hz}$ 到 $1500\mathrm{Hz}$ 连续改变时, 发现小提琴的弦只在 $880\mathrm{Hz}$ 与 $1320\mathrm{Hz}$ 这两个频率发生振动. 试问弦中的张力有多大?
\end{problem}
\clearpage


\begin{problem}{3.5.5}
一直立圆柱形管子, 上端开口, 可以往里灌水, 当管中水面离管口距离为 $15.95 \mathrm{~cm} 、 48.45 \mathrm{~cm}$ 和 $80.95 \mathrm{~cm}$ 时, 相继地观察到一个频率为 $512 \mathrm{~Hz}$ 的音叉发生的空气柱共振, 如图 3.68 所示.

(1) 求出在空气中的声速；

(2) 精确地指出该管邻近顶端波腹的位置；

(3) 如果上述测量是二年级大学生所做,你如何评价他的工作.
\begin{center}
\begin{tikzpicture}[>=Stealth, thick]
  % Three open-top vertical cylindrical tubes
  \foreach \xc in {0, 1.4, 2.8} {
    \draw (\xc-0.25, 0) -- (\xc-0.25, -4.5) -- (\xc+0.25, -4.5) -- (\xc+0.25, 0);
  }

  % Reference dashed lines and depth labels
  \draw[dashed, thin] (-0.35, 0) -- (3.85, 0) node[right] {$0$};
  \draw[dashed, thin] (-0.35, -0.80) -- (3.85, -0.80) node[right] {$15.95\,\mathrm{cm}$};
  \draw[dashed, thin] (-0.35, -2.44) -- (3.85, -2.44) node[right] {$48.45\,\mathrm{cm}$};
  \draw[dashed, thin] (-0.35, -4.08) -- (3.85, -4.08) node[right] {$80.95\,\mathrm{cm}$};

  % Tube 1 (left): water level at 80.95 cm (y = -4.08), 5/4 wavelength
  \draw[smooth, domain=0:-4.08, samples=80]
    plot ({0 + 0.19*cos(90*(0.02 - \x)/0.82)}, {\x});
  \draw[smooth, domain=0:-4.08, samples=80]
    plot ({0 - 0.19*cos(90*(0.02 - \x)/0.82)}, {\x});
  % Water surface symbol
  \draw[thin] (-0.18, -4.16) -- (0.18, -4.16);
  \draw[thin] (-0.12, -4.24) -- (0.12, -4.24);
  \draw[thin] (-0.06, -4.32) -- (0.06, -4.32);

  % Tube 2 (middle): water level at 48.45 cm (y = -2.44), 3/4 wavelength
  \draw[smooth, domain=0:-2.44, samples=60]
    plot ({1.4 + 0.19*cos(90*(0.02 - \x)/0.82)}, {\x});
  \draw[smooth, domain=0:-2.44, samples=60]
    plot ({1.4 - 0.19*cos(90*(0.02 - \x)/0.82)}, {\x});
  % Water surface symbol
  \draw[thin] (1.4-0.18, -2.52) -- (1.4+0.18, -2.52);
  \draw[thin] (1.4-0.12, -2.60) -- (1.4+0.12, -2.60);
  \draw[thin] (1.4-0.06, -2.68) -- (1.4+0.06, -2.68);

  % Tube 3 (right): water level at 15.95 cm (y = -0.80), 1/4 wavelength
  \draw[smooth, domain=0:-0.80, samples=40]
    plot ({2.8 + 0.19*cos(90*(0.02 - \x)/0.82)}, {\x});
  \draw[smooth, domain=0:-0.80, samples=40]
    plot ({2.8 - 0.19*cos(90*(0.02 - \x)/0.82)}, {\x});
  % Water surface symbol
  \draw[thin] (2.8-0.18, -0.88) -- (2.8+0.18, -0.88);
  \draw[thin] (2.8-0.12, -0.96) -- (2.8+0.12, -0.96);
  \draw[thin] (2.8-0.06, -1.04) -- (2.8+0.06, -1.04);
\end{tikzpicture}
\end{center}
\end{problem}
\clearpage


\begin{problem}{3.5.6}
小提琴上的琴弦长 L，可认为两端是固定的，空弦的基频为 $\nu_{0}$ ，提琴手在离一

端 $\frac{L}{4}$ 距离处用弓拉弦，另外用手轻轻按在弦的中点。

(1) 问在这些条件下,他能激发的最低频率多大?示意地画出弦振动时的形状;

(2) 在此条件下,一次谐波的频率多大?
\begin{center}
\begin{tikzpicture}[>=Stealth, thick]
  % Three open-top vertical cylindrical tubes
  \foreach \xc in {0, 1.4, 2.8} {
    \draw (\xc-0.25, 0) -- (\xc-0.25, -4.5) -- (\xc+0.25, -4.5) -- (\xc+0.25, 0);
  }

  % Reference dashed lines and depth labels
  \draw[dashed, thin] (-0.35, 0) -- (3.85, 0) node[right] {$0$};
  \draw[dashed, thin] (-0.35, -0.80) -- (3.85, -0.80) node[right] {$15.95\,\mathrm{cm}$};
  \draw[dashed, thin] (-0.35, -2.44) -- (3.85, -2.44) node[right] {$48.45\,\mathrm{cm}$};
  \draw[dashed, thin] (-0.35, -4.08) -- (3.85, -4.08) node[right] {$80.95\,\mathrm{cm}$};

  % Tube 1 (left): water level at 80.95 cm (y = -4.08), 5/4 wavelength
  \draw[smooth, domain=0:-4.08, samples=80]
    plot ({0 + 0.19*cos(90*(0.02 - \x)/0.82)}, {\x});
  \draw[smooth, domain=0:-4.08, samples=80]
    plot ({0 - 0.19*cos(90*(0.02 - \x)/0.82)}, {\x});
  % Water surface symbol
  \draw[thin] (-0.18, -4.16) -- (0.18, -4.16);
  \draw[thin] (-0.12, -4.24) -- (0.12, -4.24);
  \draw[thin] (-0.06, -4.32) -- (0.06, -4.32);

  % Tube 2 (middle): water level at 48.45 cm (y = -2.44), 3/4 wavelength
  \draw[smooth, domain=0:-2.44, samples=60]
    plot ({1.4 + 0.19*cos(90*(0.02 - \x)/0.82)}, {\x});
  \draw[smooth, domain=0:-2.44, samples=60]
    plot ({1.4 - 0.19*cos(90*(0.02 - \x)/0.82)}, {\x});
  % Water surface symbol
  \draw[thin] (1.4-0.18, -2.52) -- (1.4+0.18, -2.52);
  \draw[thin] (1.4-0.12, -2.60) -- (1.4+0.12, -2.60);
  \draw[thin] (1.4-0.06, -2.68) -- (1.4+0.06, -2.68);

  % Tube 3 (right): water level at 15.95 cm (y = -0.80), 1/4 wavelength
  \draw[smooth, domain=0:-0.80, samples=40]
    plot ({2.8 + 0.19*cos(90*(0.02 - \x)/0.82)}, {\x});
  \draw[smooth, domain=0:-0.80, samples=40]
    plot ({2.8 - 0.19*cos(90*(0.02 - \x)/0.82)}, {\x});
  % Water surface symbol
  \draw[thin] (2.8-0.18, -0.88) -- (2.8+0.18, -0.88);
  \draw[thin] (2.8-0.12, -0.96) -- (2.8+0.12, -0.96);
  \draw[thin] (2.8-0.06, -1.04) -- (2.8+0.06, -1.04);
\end{tikzpicture}
\end{center}
\end{problem}
\clearpage


\begin{problem}{3.5.7}
一长为 0.5m 的小提琴弦, 其基频为 200Hz.

(1) 横向脉冲在此弦上的传播速度多大?

(2) 画出一脉冲在弦的端点反射前后的图形；

(3) 画出此弦比基频较高的两种振动模式的草图, 并给出它们的频率.
\end{problem}
\clearpage


\begin{problem}{3.5.8}
在两个相距为 100cm 的固定点之间张紧一根弦，在 100Hz～350Hz 的频率范围内,仅有以下频率能被激发:160Hz、240Hz、320Hz,这些振动模式的波长各是多少?
\begin{center}
\begin{tikzpicture}[>=Stealth, thick]
  % Horizontal equilibrium string axis
  \draw (-1.4, 0) -- (3.8, 0);

  % Fixed vertical wall at x = 0 with hatching on the left
  \draw (0, -1.6) -- (0, 1.6);
  \foreach \y in {-1.45, -1.25, ..., 1.55} {
    \draw[thin] (-0.16, \y-0.14) -- (0, \y);
  }

  % Incident pulse (peak at x = +0.42): dashed for x < 0, solid for x >= 0
  \draw[dashed, smooth, domain=-0.65:0, samples=30]
    plot (\x, {1.35*exp(-(\x-0.42)^2/0.16)});
  \draw[smooth, domain=0:1.45, samples=50]
    plot (\x, {1.35*exp(-(\x-0.42)^2/0.16)});

  % Reflected inverted pulse (peak at x = -0.42): dashed for x < 0, solid for x >= 0
  \draw[dashed, smooth, domain=-1.45:0, samples=50]
    plot (\x, {-1.35*exp(-(\x+0.42)^2/0.16)});
  \draw[smooth, domain=0:0.65, samples=30]
    plot (\x, {-1.35*exp(-(\x+0.42)^2/0.16)});

  % Direction arrows and labels
  \draw[<-] (0.55, 1.35) -- (1.05, 1.35) node[right] {入射脉冲};
  \draw[->] (0.50, -0.65) -- (1.00, -0.65) node[right] {反射脉冲};
\end{tikzpicture}
\end{center}
\end{problem}
\clearpage


\begin{problem}{3.5.9}
抓住或夹住细铝杆的中点, 铝杆长 1m, 沿着杆的轴线纵向用锤子打击杆的一端, 产生 2500Hz 的声波.

(1) 求铝杆中的声速；

(2) 要激发 3750Hz 的声波, 手应抓在杆的什么地方? 打击不同端有差别吗?

（3）假如像原来那样，手握杆的中点，但在端点横向而不是纵向打击杆，定性解释为什么合成的声波频率比原来的低。
\end{problem}
\clearpage


\begin{problem}{3.5.10}
在温度为 300K 的空气中, 假设频率为 $10^{3}$ Hz 的平面驻声波引起压强变化的幅值为 $1 \, dyn/cm^{2}$ (周围压强为 $10^{6} \, dyn/cm^{2}$ ). 试估计此波动引起空气分子位移振幅的量值.
\end{problem}
\clearpage


\begin{problem}{3.5.11}
(1)对于一端封闭的管子,画出其第二振动模式的空气位移图和压强变化图;

(2) 上述第二振动模式的频率与基频有何关系?
\end{problem}
\clearpage


\begin{problem}{3.5.12}
一吉他弦线长 $80\mathrm{cm}$ , 其基频为 $400\mathrm{Hz}$ , 做此基频的振动, 使弦中点的最大位移为 $2\mathrm{cm}$ , 如弦中张力 $T_{0} = 10^{6}\mathrm{dyn}$ , 试问作用于弦支承端点的作用力的垂直于弦平衡位置的分量的最大值有多大?
\end{problem}
\clearpage


\begin{problem}{3.5.13}
一根长 L 的弦由两种线密度的弦相连而成, 右边 $\frac{2}{3}L$ 长的弦线密度为 $\eta$ , 左边 $\frac{1}{3}L$ 长的弦线密度为 $4\eta$ , 张力为 $T_{0}$ , 两端固定. 求基频和第一、第二、第三、第四泛音频率.
\end{problem}
\clearpage


\begin{problem}{3.5.14}
线密度为 $\eta$ 、长为 $l$ 的琴弦两端固定，受张力 $T_0$ 作用，设在 $t = 0$ 时，将弦的中点由平衡位置拨到距离 $s(s \ll l)$ 处，使形成一等腰三角形，然后将弦静止释放．求弦的运动．
\end{problem}
\clearpage


\begin{problem}{3.5.15}
密度为 $\rho_{1}$ 、声速为 $v_{1}$ 的介质 1，与密度为 $\rho_{2}$ 、声速为 $v_{2}$ 的介质 2 间有一不可渗透的平面界面，一列压强振幅为 A、频率为 $\nu$ 的平面声波自介质 1 垂直射向界面。

(1) 写出在界面上的边界条件；

(2) 用边界条件求出自界面反射回介质 1 的波的压强振幅 $A'$ 和透射到介质 2 的波的压强振幅 $B$ .
\end{problem}
\clearpage


\begin{problem}{3.5.16}
两根相同的钢琴弦在保持相同的张力时有相同的基频,设基频为 600Hz,试问当这两根弦同时振动时,使其中一根弦增加多大百分比的张力,每秒发生 6 个拍?
\end{problem}
\clearpage


\begin{problem}{3.5.17}
设有两列声波,一列在空气中,另一列在水中,如果这两列声波的强度相等,水中声波的压强振幅与空气中声波的压强振幅之比多大?若这两列波的压强振幅相等,它们的声强之比多大?
\end{problem}
\clearpage


\begin{problem}{3.5.18}
两端开口的长为 $l$ 的风琴管可用来测量亚音速风洞中空气的马赫数 $v / c$ ，其中 $v$ 是空气的流动速度， $c$ 是在静止空气中的声速。观察到当风琴管固定在洞中不动时，与周期为 $T$ 的基波发生共振。若 $v / c = \frac{1}{2}$ ，求 $T / T_0, T_0$ 是风琴管置于静止空气中时在基波发生共振的基波周期。
\end{problem}
\clearpage


\begin{problem}{3.5.19}
一汽笛发出频率为 $1000\mathrm{Hz}$ 的声音，该汽笛以 $10\mathrm{m / s}$ 的速率离开你向一悬崖运动，空气中的声速为 $330\mathrm{m / s}$ . 试问：

(1) 你听到的直接从汽笛传来的声波的频率多大?

(2) 你听到从悬崖反射回来的声波频率多大?

(3) 你能否听到拍频？
\end{problem}
\clearpage


\begin{problem}{3.5.20}
一卡车在其车顶上装有两个扬声器,一个朝前,一个朝后,此卡车以 $1\mathrm{m / s}$ 的速度向你驶来,如果扬声器都以 $1000\mathrm{Hz}$ 的频率振动,问你听到的由直接传来的声音和由汽车后一建筑物反射的回声形成的拍频等于多少? (取声速为 $330\mathrm{m / s}$ )
\end{problem}
\clearpage


\begin{problem}{3.5.21}
一大学生手持一振动频率为 $440\mathrm{Hz}$ 的音叉以 $1.2\mathrm{m / s}$ 的速度从一座墙走开。问从墙反射出的回声的音调高于还是低于音叉的音调？大学生听到的拍频多大？声音的速度为 $330\mathrm{m / s}$
\end{problem}
\clearpage


\begin{problem}{3.5.22}
血液在动脉中流动的速度可以用超声波多普勒频移来测量．假定频率为 $1.5\times10^{6}Hz$ 的超声波被以1m/s速度流动的血液反射回来，若超声波在人体组织中的速度为1500m/s，且超声波以非常小的角度入射，如图3.75所示．求入射和反射波之间的频率移动 $\Delta\nu$ .
\begin{center}
\begin{tikzpicture}[>=Stealth, thick]
  % Source and receiver dot + label
  \coordinate (S) at (0, 0.75);
  \fill (S) circle (1.4pt);
  \node[above left=0.12cm and -0.25cm of S, align=left, font=\small] {声源和\\接收器};

  % Concentric ultrasonic wavefront arcs
  \draw (S) ++(-35:0.30) arc (-35:30:0.30);
  \draw (S) ++(-38:0.55) arc (-38:28:0.55);
  \draw (S) ++(-40:0.82) arc (-40:26:0.82);

  % Incident beam direction arrow at small angle to artery
  \draw[->] (1.05, 0.56) -- (1.75, 0.38);

  % Artery walls
  \draw (-0.3, 0.15) -- (4.8, 0.15);
  \draw (-0.3, -0.55) -- (4.8, -0.55);

  % Blood flow velocity arrow
  \draw[->] (2.7, -0.20) -- (3.6, -0.20) node[right] {$v=1\,\mathrm{m/s}$};
\end{tikzpicture}
\end{center}
\end{problem}
\clearpage


\begin{problem}{3.5.23}
一站在铁路轨道边的学生听到火车向他驶来和离他而去的汽笛声的频率分别为 250Hz 和 200Hz, 假定声音在空气中的速度为 340m/s, 问火车的速度多大?
\end{problem}
\clearpage


\begin{problem}{3.5.24}
一个声学运动探测器发射 50kHz 的信号并接收回声, 如果回声有多普勒频移大于 100Hz, 运动物体就被记录下来. 设空气中的声速为 330m/s, 要被探测器记录, 运动物体需以多大的速率向着或离开探测器运动?
\end{problem}
\clearpage


\begin{problem}{3.5.25}
火车以 25m/s 的速率行驶, 汽笛声的频率为 500Hz, 一个人站在离铁轨 100m 处, 当人与汽笛的连线与铁轨垂直时取为 t=0, 问 t 为何值时, 人听到的汽笛声的频率比原频率低 25Hz? 空气中声速为 330m/s.
\end{problem}
\clearpage


\begin{problem}{3.5.26}
声波从波源水平地传向接收器, 波源的速率为 u, 接收器的速率为 v, 向同一方向运动, 另有速率为 w 的风也水平地从波源吹向接收器. 若波源发射的声波频率为 $\nu$ , 静止空气中的声速为 V. 设风向也与波源、接收器的运动方向一致. 试证: 接收器收到的频率为

\[
\nu = \frac {V - v + w}{V - u + w} \nu
\]

并说明此式的适用范围.
\end{problem}
\clearpage


\begin{problem}{3.5.27}
频率为 $\nu_{0}$ 的声源，在距地面 h 高度处以匀速 u 沿 x 方向做水平运动，一观察者位于地面 x=0 处，t=0 时，声源正经过他的上空，设声速为 v，试证：

(1) t 时刻接收到的信号是声源 $t_{s}$ 时发出的，

\[
t _ {s} = \frac {v ^ {2}}{v ^ {2} - u ^ {2}} \left\{t - \frac {1}{v} \left[ h ^ {2} \left(1 - \frac {u ^ {2}}{v ^ {2}}\right) + u ^ {2} t ^ {2} \right] ^ {\frac {1}{2}} \right\}
\]

(2) $t$ 时刻所接收到的信号的频率作为发出时刻 $t_s$ 的函数为

\[
\nu (t _ {s}) = \frac {\nu_ {0}}{1 + \frac {u}{v} \frac {u t _ {s}}{(h ^ {2} + u ^ {2} t _ {s} ^ {2}) ^ {1 / 2}}}
\]
\begin{center}
\begin{tikzpicture}[>=Stealth, thick]
  % Key coordinates
  \coordinate (O) at (0, 0);
  \coordinate (H) at (0, 1.6);
  \coordinate (S) at (3.2, 1.6);
  \coordinate (U) at (4.3, 1.6);

  % Observer and vertical height
  \node[below left] at (O) {$O$};
  \draw (O) -- (H) node[midway, left] {$h$};

  % Horizontal source trajectory
  \draw (H) -- (4.7, 1.6);
  \draw[->] (H) -- (0.65, 1.6) node[below] {$u$};
  \node[above left=0.04cm and 0.02cm] at (S) {$S$};

  % Line of sight OS and sound speed v
  \draw (O) -- (S);
  \draw[->] (S) -- ($(S)!0.24!(O)$) node[below right=0pt and 1pt] {$v$};

  % Velocity decomposition at S: u = u_1 + u_2
  \draw[->] (S) -- (U) node[below=2pt] {$u$};
  \coordinate (U1) at ($(S)!(U)!($(S)+(S)-(O)$)$);
  \coordinate (U2) at ($(S)+(U)-(U1)$);
  \draw[->] (S) -- (U1) node[above right] {$u_1$};
  \draw[->] (S) -- (U2) node[below right] {$u_2$};
  \draw[dashed, thin] (U1) -- (U) -- (U2);

  % Top dimension line for ut_s
  \draw[thin] (0, 2.05) -- (0, 2.35);
  \draw[thin] (3.2, 2.05) -- (3.2, 2.35);
  \draw[<->, thin] (0, 2.20) -- (3.2, 2.20) node[midway, fill=white, inner sep=2pt] {$ut_s$};

  % Bottom dimension line parallel to OS for v(t - t_s)
  \coordinate (Ooff) at ($(O)!0.55cm!-90:(S)$);
  \coordinate (Soff) at ($(S)!0.55cm!90:(O)$);
  \draw[thin] ($(O)!0.42cm!-90:(S)$) -- ($(O)!0.68cm!-90:(S)$);
  \draw[thin] ($(S)!0.42cm!90:(O)$) -- ($(S)!0.68cm!90:(O)$);
  \draw[<->, thin] (Ooff) -- (Soff) node[midway, fill=white, inner sep=1.5pt] {$v(t-t_s)$};
\end{tikzpicture}
\end{center}
\end{problem}
\clearpage


\begin{problem}{3.5.28}
频率为 500Hz 的音叉以 5.0rad/s 的匀角速在半径为 6.0m 的圆周上运动。试求在距圆心 1.2m 处的人听到的最高和最低频率各为多少？设声速为 330m/s。
\end{problem}
\clearpage


\begin{problem}{3.5.29}
声音在大气中的传播速度为 330m/s，一

过．问当观察者刚听到飞机隆隆声时，飞机离他的正上方的水平距离有多远？
\end{problem}
\clearpage


\begin{problem}{3.5.30}
若声音在空气中的速度为 c，声源沿 x 方向运动，运动速度为 v.

(1) v<c 的情况, t=0 时, 一个声脉冲由坐标原点发出, 画出 t 时刻的波阵面和 t 时刻声源位置之间的关系, 仔细地对所作的图加以标记, 以 t 时刻波源作为原点, 写出波阵面位置的方程;

(2) v>c 的情况, 声源发出连续信号, 画出由运动声源所产生的波前, 说明作图的依据, 写下联系波前形状和其他一些已知量的方程.
\end{problem}
\clearpage


\begin{problem}{3.5.31}
一架飞机在上空做水平飞行, 某时刻, 地面上的人同时听到飞机的隆隆声的位置在通过 $(2\sqrt{3} \times 10^{3}, \sqrt{2} \times 10^{2})$ 点的下述曲线上:

\[
x ^ {2} - 3 y ^ {2} - 2 \sqrt {3} \times 1 0 ^ {3} x + 6 \times 1 0 ^ {4} = 0
\]

x、y 在地面上, 均以米为单位. 能否从此曲线求出马赫角 $\alpha$ 、飞机飞行的高度 h, 并说明飞机航行的方向.
\end{problem}
\clearpage

